\documentclass[11pt]{article}
\usepackage[margin=1in]{geometry}
\usepackage{amsmath,amssymb,amsthm,mathtools}
\usepackage{hyperref}
\usepackage{enumitem}
\usepackage{booktabs}
\usepackage{longtable}
\usepackage{array}
\usepackage{microtype}
\setlength{\emergencystretch}{2em}
\hypersetup{colorlinks=true,linkcolor=blue,urlcolor=blue,
pdftitle={The Cobordism Functor, the Theta Register, and Root Systems}}

\newtheorem{theorem}{Theorem}[section]
\newtheorem{lemma}[theorem]{Lemma}
\newtheorem{proposition}[theorem]{Proposition}
\newtheorem{corollary}[theorem]{Corollary}
\theoremstyle{definition}
\newtheorem{definition}[theorem]{Definition}
\newtheorem{example}[theorem]{Example}
\theoremstyle{remark}
\newtheorem{remark}[theorem]{Remark}

\newcommand{\Z}{\mathbb{Z}}
\newcommand{\R}{\mathbb{R}}
\newcommand{\C}{\mathbb{C}}
\newcommand{\im}{\operatorname{im}}
\newcommand{\rank}{\operatorname{rank}}
\newcommand{\Hgroup}{\mathcal{H}}
\newcommand{\Sp}{\operatorname{Sp}}
\newcommand{\GL}{\operatorname{GL}}
\newcommand{\SL}{\operatorname{SL}}
\newcommand{\Tr}{\operatorname{Tr}}
\newcommand{\ann}{\operatorname{ann}}
\DeclareMathOperator{\Span}{span}
\DeclareMathOperator{\Hom}{Hom}

\title{The Cobordism Functor, the Theta Register, and Root Systems\\[4pt]
\large Consolidated derivation and experiment specification}
\author{}
\date{Consolidated and revised 14 September 2026}

\begin{document}
\maketitle

\begin{abstract}
This is the single authoritative account of the harmonic-operator
experiment: its analytical derivation, geometric operator algebra, theta
tensor register, and root-system extension. The current objective is an
exact representation of finite-dimensional algebra by quantized geometry,
not synthesis of every operator as one engine-generated spacetime.
Harmonic forms parameterize a classical phase space; quantization supplies
the quantum Hilbert space. Geometric translations and their complex linear
combinations form its full matrix algebra, preserving tensor products,
adjoints, states and operator vectors. Root-lattice geometry represents
characters and fusion with explicitly declared quantization data.
Single-cobordism synthesis and physical dynamics remain separate questions.
The text distinguishes theorems, source-reported measurements and open
constructions. Companion problems P1--P46 and their solutions remain study
aids, not competing specifications.
\end{abstract}

\begingroup\small
\section*{Authority and provenance}
\addcontentsline{toc}{section}{Authority and provenance}
Edit \texttt{docs/design/cobordism\_functor\_reading.tex} in the Tessera
repository; the adjacent PDF is generated from that source. This is the
authoritative copy, included in Git at the owner's request. The files in
\texttt{\string~/tessera-notes/cobordism-functor/} are synchronized reading
copies, not a second editing authority. Regenerate the PDF and synchronize
those copies after an accepted source revision. This consolidation incorporates the original chat derivation
of 13 September 2026, the previous reading (including its root-system
development), and \texttt{docs/design/qubit\_cobordism\_spec.md} at Tessera
revision
\href{https://github.com/akellehe/tessera/blob/3d54e5990240bc68b0c4774b7a9386dd2c6c6d37/docs/design/qubit_cobordism_spec.md}{\texttt{3d54e5990240}} (full revision encoded in this link).
It supersedes their conflicting claims. The former
\texttt{qubit\_cobordism\_spec.md} was retired; its replacement location
notice now points to this maintained repository source. This revision changes the research contract and
tickets, not the implementation. Earlier notes are preserved in
\texttt{archive/pre-consolidation-2026-09-14/};
\texttt{original\_derivation.md} is historical, not normative.
Reported numerical results below are inherited from the pinned sources,
not measurements rerun during consolidation, except
Section~\ref{sec:evidence-merged}, which records measurements made on 14
September 2026 after consolidation (Tessera pull requests \#1116, \#1119
and \#1120, merged) and the owner's decisions on what is measured next
(Section~\ref{sec:next}). The subsequent owner-approved revision makes
representation, rather than native synthesis, the acceptance goal.
Historical numbers retain their original conventions; they are not
certificates for corrected claims. Companion maps, problems and solutions
are non-normative study aids and may use older conventions.

\paragraph{Notation.}
$D$ is geometric dimension, $p$ cochain degree, $r_H$ harmonic dimension,
$g$ the dimension of a polarized Jacobian, and $r$ the rank of a Lie algebra.
The primer retains $k$ for cochain degree in its Hodge discussion and for
positive integer quantization level in its theta/affine discussion; those
are independent quantities. In $\mathfrak{su}(N)$, $N$ is the defining
representation dimension, not simplex dimension. $M$ denotes period
transport; $\mathsf M$ denotes its standard Siegel-coordinate counterpart
where that bridge applies; $M_p$ denotes a mass matrix. $\Omega$ is an intersection matrix
in the restriction discussion and a complex period matrix in the theta
discussion, as explicitly indicated. Transpose and Hermitian adjoint are
never interchangeable without the stated real or duality hypotheses.

\endgroup
\clearpage
\tableofcontents
\clearpage

\part{Primer}

\section{What is being claimed}

Let $W$ be a compact oriented triangulated $3$-manifold whose boundary is two tori,
$\partial W = M_0 \sqcup M_1$. Each torus carries a \emph{marking}: two
cycles $A, B$ that form a basis of $H_1(M_i;\Z) \cong \Z^2$; their periods
give the dual cohomology coordinates. The claim has
three levels. The primary goal is an equivalent geometric representation
of algebra, including complex linear combinations of geometric operations.
Bare geometry is not claimed to choose a state, level or physical dynamics.

\begin{enumerate}[nosep]
  \item (\emph{Exact.}) The degree-one harmonic space of $W$, restricted to
    the boundary, is a subspace $L_W \subset H^1(M_0)\oplus H^1(M_1)$ that is
    Lagrangian for the boundary intersection form and, when it is the graph
    of a map, defines a metric-independent matrix $M$. If the restrictions
    identify the full integral lattices, $M$ is integral and unimodular.
    Gluing cobordisms multiplies their matrices. $M$ does not depend on the
    metric.
  \item (\emph{Exact, on a new register.}) Quantizing each marking lattice at
    level $k$ gives a $k$-dimensional space $\Theta_k$ per torus; a disjoint
    union of tori gives the tensor product; an appropriately lifted symplectic
    period map acts by the unitary projective matrix $\rho_k(\mathsf M)$
    after the coordinate bridge, with a declared
    polarization, theta structure and lift convention. Geometric clock and
    shift operations span the full operator algebra
    (Theorem~\ref{thm:full-algebra}), not only the modular gates. At even level,
    including $k=2$, these operators are Clifford; antisymplectic maps
    require an antiunitary interpretation.
  \item (\emph{Exact representation with declared root data.}) Replacing
    the abelian lattice by a root system
    gives the non-abelian register: Weyl-antisymmetrized theta functions,
    Kac--Peterson modular matrices, Verlinde fusion, and --- with punctures
    --- braid representations with proven density in specified sectors.
    An internal scalar Euclidean lattice realizes $A_2$ explicitly
    (Proposition~\ref{prop:scalar-a2}). Level and affine reduction are
    declared structure. Native bulk synthesis, physical sewing and braid
    synthesis are later programs, not prerequisites for this representation.
\end{enumerate}

For fixed topology, markings and flat local system, the metric controls
harmonic representatives, pairings and the validity of the harmonic model.
It also affects metric-dependent observations and dynamics. None of this
makes the period relation itself metric-dependent.

\begin{definition}[Representation versus synthesis]\label{def:representation}
A geometric representation is a geometric vector space or algebra with an
invertible map to its algebraic counterpart preserving specified operations:
here addition, complex scalar multiplication, composition, adjoint and
tensor product. Symbols may be complex linear combinations of decorated
cycles; their coefficients carry information and need not emerge from a
metric. Synthesis asks additionally for one allowed scalar simplex complex
whose native readout realizes a specified operator. The current experiment
requires the first statement, not the second.
\end{definition}

\section{Cochains, the Laplacian, and the harmonic space}\label{sec:cochains}

A $k$-\emph{chain} is a formal integer combination of oriented
$k$-simplices; $\partial_k$ is the boundary map, $\partial_{k-1}\partial_k = 0$.
A $k$-\emph{cochain} is a linear function on $k$-chains: one number per
oriented simplex, changing sign with orientation. The \emph{coboundary}
$d_k = \partial_{k+1}^{\mathsf T}$ satisfies $\langle d\omega, c\rangle =
\langle \omega, \partial c\rangle$ and $d_{k+1} d_k = 0$. The \emph{period}
of a cochain over a cycle $z$ ($\partial z = 0$) is $\langle \omega, z\rangle$.

\begin{lemma}\label{lem:exact}
A coboundary has zero periods: $\langle d\eta, z\rangle = \langle \eta,
\partial z\rangle = 0$ for every cycle $z$.
\end{lemma}

Cohomology $H^k = \ker d_k / \im d_{k-1}$ is the quotient of closed cochains
by coboundaries; its dimension is the Betti number $b_k$.

A metric enters through positive matrices $M_k$ on $k$-cochains (for us the
Whitney mass, a function of the edge lengths) and gives the codifferential
$\delta_k = M_{k-1}^{-1} d_{k-1}^{\dagger} M_k$ and the Hodge Laplacian
$L_k = d_{k-1}\delta_k + \delta_{k+1} d_k$. The \emph{harmonic space} is
$\Hgroup_k = \ker L_k = \ker d_k \cap \ker \delta_k$.

\begin{theorem}[Discrete Hodge]\label{thm:hodge}
If the cochain masses are positive Hermitian then $C^k = \im d_{k-1} \oplus \Hgroup_k \oplus
\im \delta_{k+1}$, orthogonally, and $\Hgroup_k \to H^k$ (a harmonic cochain
to its class) is an isomorphism. In particular $\dim \Hgroup_k = b_k$ and
every class has exactly one harmonic representative.
\end{theorem}

\begin{remark}[The Lorentzian caveat]
On complex lengths $M_k$ is complex bilinear and not positive, and
the positive Hodge theorem does not apply. Equality of dimensions alone
does not imply a harmonic identification with cohomology. Certify closedness,
coclosedness, an independent kernel of the same span, and that the map to
cohomology is injective and onto (for example by a full period-rank test).
The code's H1--H3 checks form part of this contract; see
Section~\ref{sec:kernel-contract}. All uses of harmonic representatives in
the cohomological theorems assume this identification. Quantum statements
add a positive Hermitian Hilbert metric explicitly.
\end{remark}

\begin{theorem}[Periods do not move]\label{thm:invariance}
Fix a class $\alpha \in H^k$ and a cycle $z$. If $\omega$ and $\omega'$ are
the harmonic representatives of $\alpha$ for two metrics, then
$\langle \omega, z\rangle = \langle \omega', z\rangle$.
\end{theorem}
\begin{proof}
Both represent $\alpha$, so $\omega - \omega' = d\eta$; apply
Lemma~\ref{lem:exact}.
\end{proof}

\section{Restriction to the boundary}\label{sec:restriction}

Write $\iota: \partial W \hookrightarrow W$ for the inclusion and
$\iota^*: H^k(W) \to H^k(\partial W)$ for restriction of classes. On
cochains this restricts values to boundary cells. A closed cochain restricts
to a closed cochain, but generally not to a coclosed one for the boundary's
own metric. Passing to classes, or then taking the boundary harmonic
representative when defined, is essential. Let $m = \dim \partial W$
($m = 2$ for us).

\begin{theorem}[Mutual annihilators]\label{thm:annihilators}
Under the Poincar\'e pairing $H^k(\partial W) \times H^{m-k}(\partial W) \to \R$,
\[
  \im\big(H^k(W) \to H^k(\partial W)\big)
  \;=\; \ann\Big(\im\big(H^{m-k}(W) \to H^{m-k}(\partial W)\big)\Big).
\]
\end{theorem}

The proof has two halves. Orthogonality is Stokes with the cup product: for
closed $\alpha \in C^k(W)$ and $\beta \in C^{m-k}(W)$,
\[
  \langle \alpha|_{\partial W} \cup \beta|_{\partial W}, [\partial W]\rangle
  = \langle \alpha \cup \beta, \partial [W]\rangle
  = \langle d(\alpha\cup\beta), [W]\rangle = 0,
\]
since $d(\alpha\cup\beta) = d\alpha \cup \beta \pm \alpha \cup d\beta = 0$.
That the two images are \emph{exact} annihilators, not merely orthogonal, is
a dimension count from the long exact sequence of the pair $(W, \partial W)$
together with Lefschetz duality $H^k(W,\partial W) \cong H_{m+1-k}(W)$. Both
halves hold for compact oriented triangulated manifolds over a field;
arbitrary simplicial complexes need not satisfy Poincar\'e--Lefschetz
duality. No metric or signature appears.

\begin{corollary}[Half lives, half dies]\label{cor:half}
For $\partial W$ a union of tori and $k = 1$: $\rank(H^1(W) \to H^1(\partial W))
= \tfrac12 b_1(\partial W)$, and the image is a Lagrangian subspace for the
intersection form on $H^1(\partial W)$.
\end{corollary}

For two tori, $b_1(\partial W) = 4$ and exactly two of the four boundary
periods can be prescribed; the bulk decides which two. Call the image
$L_W$. It is a linear \emph{relation} between $H^1(M_0)$ and $H^1(M_1)$.

\begin{remark}[The orientation sign]\label{rem:sign}
Write the intersection form on one torus in its marking as $J =
\begin{psmallmatrix} 0 & 1\\ -1 & 0\end{psmallmatrix}$. The boundary form is
$\Omega = J_A \oplus \varepsilon J_B$ where $\varepsilon = \pm 1$ records that
the two ends inherit \emph{opposite} orientations from the bulk, corrected
by the orientation of any relabelling $\varphi$ applied to the far surface:
$\varepsilon = -\det\varphi^*$. This sign is declared data, never searched
for. Problem~P4 shows why it cannot be read off $\det M$ once
there are four tori. The displayed rule applies to the declared collar
markings. With more components, assemble one signed block per boundary
component from its actual induced orientation; a total determinant is
insufficient. Orientation is needed for this pairing even when it is not
used to define charge.
\end{remark}

\section{The monodromy}\label{sec:monodromy}

Let $Z$ be a basis of $\Hgroup_1(W)$ and let $P_A, P_B$ be the matrices of
periods of $Z$ over the two markings ($2 \times r_H$ each). When $P_A$ is
invertible, $L_W$ is the graph of the map
\[
  M \;=\; P_B P_A^{-1}, \qquad p_{\text{out}} = M\, p_{\text{in}}.
\]

\begin{proposition}\label{prop:basisfree}
$M$ does not depend on the basis $Z$: a change $Z \mapsto ZS$ sends
$P_A \mapsto P_A S$, $P_B \mapsto P_B S$, and $P_B S (P_A S)^{-1} = P_B P_A^{-1}$.
\end{proposition}

\begin{theorem}[Integrality under integral lattice identification]\label{thm:integer}
If the two restrictions identify the bulk lattice (modulo invisible classes
and torsion) with each full boundary integral lattice, then
$M \in \GL(2,\Z)$: it is an integer matrix with $\det M = \pm 1$.
\end{theorem}
\begin{proof}
$M$ is, by this additional hypothesis, the matrix of a $\Z$-module
isomorphism between the two rank-two lattices, hence integer with a unit
determinant.
\end{proof}

\begin{corollary}
No change of metric changes $M$: it was defined without one.
\end{corollary}

An invertible $P_A$ over $\C$ alone does not prove the integral hypothesis.
For example, $P_A=\operatorname{diag}(1,2)$ and
$P_B=\operatorname{diag}(2,1)$ give $M=\operatorname{diag}(2,1/2)$, a
rational symplectic matrix, not an integral one. This is an algebraic
counterexample to the inference, not a manifold constructed here.
Mapping cylinders satisfy the integral-isomorphism hypothesis.

For a product collar $T^2 \times [0,1]$, $M = I$. Relabelling the far end
by a simplicial automorphism $\varphi$ of the torus gives $M = \varphi^*$.
On the square-grid triangulation with one diagonal an orientation-
reversing such $\varphi$ is the swap $(i,j)\mapsto(j,i)$, with $\varphi^* =
\begin{psmallmatrix}0&1\\1&0\end{psmallmatrix}$, $\det = -1$.

\section{From a graph relation to an operator vector}\label{sec:choi}

For finite-dimensional $A, B$, the Choi--Jamio\l{}kowski identification is
$\Hom(A,B) \cong A^*\otimes B$: a map $T$ corresponds to the vector
\[
  |T\rangle\!\rangle \;=\; \sum_i e_i^* \otimes T e_i ,
\]
and its graph $\{(x, Tx)\} \subset A \oplus B$ is a subspace of dimension
$\dim A$. Conversely a subspace of $A\oplus B$ that meets $0 \oplus B$
trivially and projects onto $A$ is the graph of exactly one map.

Thus $L_W\leftrightarrow M\leftrightarrow |M\rangle\!\rangle$ is an exact
correspondence \emph{when the graph conditions hold}; it is not literal
equality of a harmonic space, a graph subspace and a vector. The dual input
factor must not be dropped. In Hilbert notation it becomes a conjugate
input factor, with ordering fixed in Section~\ref{sec:quantum-contract}.
A nonzero output with zero input violates uniqueness; an input which cannot
extend violates coverage. A map taking a nonzero input to zero is still a
valid graph and need not be invertible. No physical creation or annihilation
interpretation follows from these linear-algebra facts alone.

\section{Composition: gluing is matrix multiplication}\label{sec:composition}

Let $W = W_1 \cup_{M_1} W_2$ with $W_1: M_0 \to M_1$ and $W_2: M_1 \to M_2$.

\begin{theorem}\label{thm:compose}
$L_W = L_{W_2} \circ L_{W_1}$ as relations; when both are graphs,
$M(W) = M(W_2)\,M(W_1)$. The cylinder $M\times[0,1]$ gives $M = I$, and a
disjoint union gives the direct sum.
\end{theorem}
\begin{proof}[Proof sketch]
Mayer--Vietoris for $W = W_1 \cup W_2$ with $W_1 \cap W_2 = M_1$: a
cohomology class on $W$ restricts to classes on $W_1$ and $W_2$ that
agree in $H^1(M_1)$, and exactness says each such compatible pair extends.
The extension need not be unique: the connecting map from $H^0(M_1)$
can contribute invisible classes. Restricted harmonic representatives
need not themselves be harmonic. Compatibility of classes gives exactly
the composition of external boundary relations. (On a simplicial complex the
glued object must actually be a complex --- one cell structure on the seam
--- or the sequence does not apply.)
\end{proof}

This gives a functor to linear relations, with disjoint union mapped to
direct sum; on the graph subcategory it gives linear maps. It is not by
itself an Atiyah vector-space TQFT, whose monoidal operation is tensor
product. Theta factorization supplies tensor spaces on objects, but a full
quantization of every cobordism relation requires further construction.

\section{What the metric controls}\label{sec:metric}

For the cohomological correspondence the relevant distinctions are:

\begin{enumerate}[nosep]
  \item Whether the harmonic space actually identifies with cohomology
    --- automatic for positive Hodge metrics, a certificate otherwise.
  \item The boundary pairings, hence whether $M$ preserves them. Green's
    formula gives the \emph{annihilator} pairing (tangential against normal),
    not an isometry: nothing in Sections~\ref{sec:restriction}--\ref{sec:composition}
    makes $M$ unitary. Dual/primal bilinear transport and positive Hermitian
    quantum isometry require distinct tests and explicit coordinate conventions
    (Section~\ref{sec:gram-contract}).
  \item Which cochain in each class is harmonic, and the frame it is read
    in. Theorem~\ref{thm:invariance}: the representative moves, the periods
    do not.
\end{enumerate}

\begin{remark}[Gauge covariance of the pairing]\label{rem:gauge}
For scalar local system $U$, use its dual $U^{-1}$ in the bilinear pairing.
If primal and dual coordinates transform by $\rho$ and $\tilde\rho$, then
the dressed mass transforms as
$G' = \tilde\rho^{-\mathsf T}G\rho^{-1}$, so
$(\tilde\rho\tilde c)^{\mathsf T}G'(\rho c)=\tilde c^{\mathsf T}Gc$.
Pairing two primal frames is not a substitute. An own-metric biorthogonal
frame can satisfy $\tilde\Phi^{\mathsf T}M_{\rm own}\Phi=I$; that does not
make a whole-metric observation an orthogonal projection, nor make its
period-coordinate Gram the identity. See Section~\ref{sec:gram-contract}.
\end{remark}

\begin{proposition}[Pinned-boundary defect is invariant]\label{prop:pinned-defect}
Fix own-boundary pairings $G_A,G_B$, coordinate markings and primal/dual
period transports $M,M^\vee$. Then
\[
 D=\frac{\|(M^\vee)^{\mathsf T}G_BM-G_A\|_F}{\|G_A\|_F}
\]
is independent of interior metric changes, including relaxation.
\end{proposition}
\begin{proof}
Every argument of the displayed function is fixed. Own-boundary Grams
depend on boundary data, not the bulk metric. At zero connection primal
and dual transports coincide; otherwise this requires proof. This tests
invariance, not restoration of positive quantum unitarity.
\end{proof}

\section{Weights add, characters multiply}\label{sec:weights}

For a torus group $T = \mathfrak h/\Lambda$ a representation decomposes into
characters $e^{2\pi i\langle\mu,\cdot\rangle}$ labelled by \emph{weights}
$\mu \in \Lambda^*$. Characters multiply under tensor product; their weights
\emph{add}: the weights of $V\otimes V'$ are the sums of a weight of $V$ and a
weight of $V'$. And the weight lattice of a product torus $T_1 \times T_2$
is $\Lambda_1^* \oplus \Lambda_2^*$. So

\[
  \text{direct sum of weight lattices} \;\longleftrightarrow\;
  \text{tensor product of representations}.
\]

This motivates, but does not construct, the finite theta register.
An integral symplectic lattice, a positive polarization and a level line
bundle define $H^0(\operatorname{Jac}(M),L^k)$; all lattice characters alone
would not give this finite-dimensional Hilbert space. Products of polarized
Jacobians with external-product line bundles give tensor products of section
spaces. Lie algebras (Sections~\ref{sec:lie}--\ref{sec:affine}) handle the
non-abelian case, where adding weights must be followed by reflecting back
into a fundamental domain of the Weyl group (Proposition~\ref{prop:klimyk}).

\section{Harmonic phase space and scalar quantization}\label{sec:phase-space}

\begin{definition}[Flat phase configuration and Jacobian]
On a closed oriented genus-$g$ surface $\Sigma$, let $\phi$ be a real
ordinary closed $1$-cochain. Gauge transformations add $d\chi$; large gauge
transformations shift its periods by $2\pi\Z$. Define
\[
 a_i=\frac1{2\pi}\oint_{A_i}\phi,\qquad
 b_i=\frac1{2\pi}\oint_{B_i}\phi,\qquad
 J(\Sigma)=H^1(\Sigma;\R)/H^1(\Sigma;\Z).
\]
The Wilson loops are $e^{2\pi i a_i},e^{2\pi i b_i}$; a real lift is
coordinate data, not another observable. Positive Hodge masses select a
unique harmonic representative of each real cohomology class. The
integral lattice and intersection pairing give the principal polarization,
represented by $2\pi\sum_i da_i\wedge db_i$.
\end{definition}
For symmetric $\Omega$ with $\operatorname{Im}\Omega>0$, use
\[
 z=b-\Omega a\pmod{\Z^g+\Omega\Z^g}.
\]
This fixes the real-holonomy to complex-Jacobian coordinate map. Positive
cochain masses prove Hodge decomposition but do not alone prove that a
discretized star yields a Riemann period matrix; certify the complex
structure or establish a controlled approximation
(Section~\ref{sec:validation}).

\paragraph{Configuration versus charged differential.}
Here $\phi$ is a configuration variable and $d$ remains the ordinary
coboundary. It is not inserted into $d_U$ to define a charged harmonic
register. On a torus with holonomies $u,v$, the minimal cellular
differentials of the rank-one local system are
\[
 d_0=\begin{pmatrix}u-1\\v-1\end{pmatrix},\qquad
 d_1=\begin{pmatrix}-(v-1)&u-1\end{pmatrix}.
\]
For $(u,v)\ne(1,1)$ both have rank one and all cohomology groups vanish.
Thus twisting the pencil with the state-selection phases changes the
problem. Store/read phase configurations separately from neutral Hodge
assembly; reuse scalar Wilson-loop transport without mutating that assembly.

\begin{definition}[Quantized register]
Choose an integer $k>0$, a positive polarization and a theta structure
(a compatible lift of the finite translations to a line bundle). The
register is $\mathcal Q_k(\Sigma)=H^0(J(\Sigma),L^k)$ with its positive
section norm, where $L^k$ is the level-$k$ line bundle. It is a quantization
of harmonic phase space, not equality with $\ker\Delta_1$. Its line
connection is scalar and lives over the Jacobian, with curvature encoding
the polarization; it is not a flat matrix-valued connection on spacetime.
\end{definition}

\section{Theta functions: the register on one torus}\label{sec:theta}

Let $\tau$ lie in the upper half-plane and $\Lambda = \Z + \tau\Z$. The
Jacobian $\C/\Lambda$ is the torus itself, and a marking $(A,B) = (1,\tau)$.

\begin{definition}[Level-$k$ theta functions]\label{def:theta}
For $j \in \Z/k$,
\[
  \theta_j(z;\tau) \;=\; \sum_{n\in\Z}
  \exp\!\Big(\pi i k\big(n+\tfrac{j}{k}\big)^2\tau + 2\pi i k\big(n+\tfrac{j}{k}\big) z\Big).
\]
$\Theta_k(\tau) = \Span\{\theta_0,\dots,\theta_{k-1}\}$ is the register's
Hilbert space for one torus; $k = 2$ is a qubit.
\end{definition}

The sum converges absolutely because $\operatorname{Im}\tau > 0$. Each
$\theta_j$ is periodic under $z \mapsto z+1$ and \emph{quasi}-periodic under
$z \mapsto z+\tau$:
\[
  \theta_j(z+\tau;\tau) = e^{-\pi i k \tau - 2\pi i k z}\,\theta_j(z;\tau),
\]
which makes $\theta_j$ a section of a line bundle of degree $k$ over the
torus rather than a function. The quantity
$\|\theta_j\|(z) = |\theta_j(z;\tau)|\,e^{-\pi k (\operatorname{Im} z)^2/\operatorname{Im}\tau}$
\emph{is} doubly periodic (Problem~P13), and each $\theta_j$ has
exactly $k$ zeros in a fundamental domain (Problem~P14).

Translations by the $k$-torsion points act on the basis:
\[
  \theta_j\big(z+\tfrac1k\big) = e^{2\pi i j/k}\theta_j(z), \qquad
  \theta_j\big(z+\tfrac{\tau}{k}\big) = e^{-\pi i \tau/k - 2\pi i z}\,\theta_{j+1}(z),
\]
the \emph{clock} and \emph{shift} operators $\mathsf Z, \mathsf X$ of a
$k$-level system, with $\mathsf Z\mathsf X = e^{2\pi i/k}\mathsf X\mathsf Z$.
They generate the Heisenberg--Weyl group of the torus's $k$-torsion.
These are operations on sections. Matrices on a column of evaluated basis
functions and on section coefficients are transposes of one another;
coherent-state transport is fixed in Section~\ref{sec:coherent}.

\section{Direct sum to tensor product}\label{sec:tensor}

For $g$ tori take a period matrix $\Omega$: $g\times g$, symmetric, with
positive-definite imaginary part (the Siegel upper half-space). For a
disjoint union $\Omega = \operatorname{diag}(\tau_1,\dots,\tau_g)$. With
characteristics $j \in (\Z/k)^g$,
\[
  \theta_j(z;\Omega) = \sum_{n\in\Z^g}
  \exp\!\Big(\pi i k\big(n+\tfrac{j}{k}\big)^{\!\mathsf T}\Omega\big(n+\tfrac{j}{k}\big)
  + 2\pi i k\big(n+\tfrac{j}{k}\big)^{\!\mathsf T} z\Big).
\]

\begin{proposition}\label{prop:factor}
If $\Omega = \operatorname{diag}(\tau_1,\tau_2)$ then $\theta_{(j_1,j_2)}(z;\Omega)
= \theta_{j_1}(z_1;\tau_1)\,\theta_{j_2}(z_2;\tau_2)$, so
$\Theta_k(\tau_1\oplus\tau_2) = \Theta_k(\tau_1)\otimes\Theta_k(\tau_2)$.
\end{proposition}
\begin{proof}
The quadratic form and the linear term both split, so the sum over $\Z^2$
is a product of sums over $\Z$.
\end{proof}

\begin{theorem}[Dimension and positive section metric]\label{thm:theta-dimension}
For a principally polarized $g$-dimensional complex torus and $k>0$,
the theta functions indexed by $j\in(\Z/k)^g$ form a basis of
$\mathcal Q_k$, of dimension $d=k^g$. Put $Y=\operatorname{Im}\Omega$ and
$z=x+\Omega y$. The positive Hermitian norm making this basis orthonormal is
\[
 \langle f,h\rangle=(2k)^{g/2}\sqrt{\det Y}
 \int_{[0,1]^{2g}}\overline{f(x+\Omega y)}h(x+\Omega y)
 e^{-2\pi k y^{\mathsf T}Yy}\,dx\,dy.
\]
\end{theorem}
\begin{proof}
Periodicity in $x$ gives a Fourier series. The $\Omega$ quasi-periodicity
determines all coefficients from their $k^g$ residue classes modulo $k$;
the theta series realize those independent choices. Integration in $x$
makes distinct frequencies orthogonal. Summing the translates of the $y$
cube gives the Gaussian integral
$(2k)^{-g/2}(\det Y)^{-1/2}$, canceled by the stated normalization.
\end{proof}

\section{The full geometric operator algebra}\label{sec:operator-algebra}

\begin{definition}[Reduced geometric translation algebra]
Let $\omega=e^{2\pi i/k}$ and $p,q\in(\Z/k)^g$. The complex vector space
$\mathcal A_{g,k}$ has basis $e_{p,q}$, representing ordered translations
around the two marked cycle families with a fixed scalar lift. Set
\[
 e_{p,q}e_{p',q'}=\omega^{q\cdot p'}e_{p+p',q+q'},\qquad
 e_{p,q}^*=\omega^{p\cdot q}e_{-p,-q}.
\]
Extend the product bilinearly and the involution conjugate-linearly;
$e_{0,0}$ is the identity. The commutator phase is the exponential of
the signed intersection of the translation labels. These scalar crossing
phases, not ordinary commutative multiplication of functions, define the
quantized geometric multiplication.
\end{definition}

\begin{theorem}[Faithful and complete geometric matrix algebra]\label{thm:full-algebra}
In the orthonormal theta basis $|j\rangle$, put
\[
 X_r|j\rangle=|j+\epsilon_r\rangle,\qquad
 Z_r|j\rangle=\omega^{j_r}|j\rangle,\qquad D_{p,q}=X^pZ^q.
\]
Then $\mathfrak Q(e_{p,q})=D_{p,q}$ is a unital $*$-algebra isomorphism
\[
 \mathcal A_{g,k}\cong\operatorname{End}(\mathcal Q_k),\qquad
 T=\frac1d\sum_{p,q}\Tr(D_{p,q}^\dagger T)D_{p,q}.
\]
Every matrix has a unique reduced geometric symbol, preserving addition,
composition and adjoint.
\end{theorem}
\begin{proof}
Theta translations give $Z_rX_s=\omega^{\delta_{rs}}X_sZ_r$, which proves
the product and involution laws. A nonzero shift has zero trace; a
nonconstant clock character sums to zero. Thus
\[
 \Tr(D_{p,q}^\dagger D_{p',q'})=d\,\delta_{p,p'}\delta_{q,q'}.
\]
The $k^{2g}=d^2$ images form an orthogonal basis of the matrix algebra,
proving injectivity, surjectivity and the inverse formula. Associativity
also follows, or can be checked from the scalar cocycle.
\end{proof}

\begin{corollary}[States, observables and Born probabilities]\label{cor:geometric-states}
Define $\operatorname{tr}_{G}(a)=d\,a_{0,0}$. Density operators correspond
to positive symbols of trace one; positivity means $a=b^*b$ for some $b$.
Pure state rays correspond to $a=a^*=a^2$ with trace one. Self-adjoint
symbols are observables. Expectations and measurement probabilities are
$\operatorname{tr}_{G}(ao)$ and $\operatorname{tr}_{G}(ae)$, for an effect
$0\le e\le1$.
\end{corollary}
\begin{proof}
The isomorphism preserves adjoints, products and trace. Positive matrices
have square roots; trace-one projections have rank one. Transport these
facts through $\mathfrak Q$.
\end{proof}

Geometric multiplication can be evaluated without matrix multiplication:
\[
 (a\star b)_{r,s}=\sum_{p,q}a_{p,q}b_{r-p,s-q}
                 \omega^{q\cdot(r-p)}.
\]
This finite twisted convolution is the computational operation to
validate; matrices are independent oracles. The coefficient count remains
$d^2$: no compression or speedup is asserted. Linear combinations of
geometric symbols need not be single curves or single engine spacetimes.

\begin{theorem}[Tensor and operator-vector compatibility]\label{thm:geometric-tensor}
For declared factor markings and theta structures,
\[
 \mathcal A_{g_1+g_2,k}\cong\mathcal A_{g_1,k}\otimes\mathcal A_{g_2,k},
 \qquad D_{(p_1,p_2),(q_1,q_2)}=D_{p_1,q_1}\otimes D_{p_2,q_2}.
\]
Output-major vectorization takes operators to
$\mathcal Q_{\rm out}\otimes\overline{\mathcal Q_{\rm in}}$.
The $d^2$ vectors $|D_{p,q}\rangle\!\rangle/\sqrt d$ form an orthonormal
operator-vector basis. Contraction of the dual intermediate factors of
$|T_2\rangle\!\rangle\otimes|T_1\rangle\!\rangle$ yields
$|T_2T_1\rangle\!\rangle$.
\end{theorem}
\begin{proof}
Cycle labels and crossing exponents split across factors; theta
factorization supplies the Hilbert-space identification. Vectorization
has inner product $\Tr(T^\dagger S)$. Contracting intermediate indices
gives $\sum_j(T_2)_{ij}(T_1)_{j\ell}$. Rectangular operators use
matrix-unit symbols between distinct registers, not a square Weyl basis
of the wrong dimension.
\end{proof}
The contraction identity uses unnormalized operator vectors. If vectors
are normalized as states, retain their norm factors; separately
renormalizing each intermediate contraction would discard amplitude and
success-probability information.
For a connected surface the tensor labels are a declared register
identification, not a claim of disjoint union. A coupled period matrix
need not factor its theta functions. The decorated-cylinder interpretation
uses linear combinations of curves modulo crossing/lift relations.
This establishes composition for this geometric algebra; quantizing
every bare three-cobordism additionally requires TQFT gluing data.
Cohomological restriction alone is not that construction.

\section{The modular action and the Weil representation}\label{sec:weil}

\begin{proposition}[Period-to-modular coordinate bridge]\label{prop:coordinate-bridge}
Let an orientation-preserving period map act on $(a,b)$ by
$M_{\rm per}=\begin{psmallmatrix}P&Q\\R&S\end{psmallmatrix}$, and suppose
the output polarization is the corresponding change of marking. Then
\[
 \Omega'=(R+S\Omega)(P+Q\Omega)^{-1},\qquad
 z'=(P+Q\Omega)^{-\mathsf T}z .
\]
Its standard Siegel matrix is
$\mathsf M=\begin{psmallmatrix}S&R\\Q&P\end{psmallmatrix}$.
\end{proposition}
\begin{proof}
The old holomorphic period blocks $(I,\Omega)$ become
$(P+Q\Omega,R+S\Omega)$; renormalize the first block.
Substituting $a'=Pa+Qb$, $b'=Ra+Sb$ in $b'-\Omega'a'$ and using
symplecticity gives $(S-\Omega'Q)(b-\Omega a)$, with
$S-\Omega'Q=(P+Q\Omega)^{-\mathsf T}$.
\end{proof}
This is a change-of-marking theorem, not a claim that an arbitrary
cobordism's output metric is the transported input metric. In genus one,
$\tau'=(M_{21}+M_{22}\tau)/(M_{11}+M_{12}\tau)$.
Never feed raw period $M_{\rm per}$ into a standard Siegel formula without
this coordinate bridge. Below $\mathsf M$ always denotes the Siegel matrix.

A symplectic integer matrix $\mathsf M = \begin{psmallmatrix}A&B\\C&D\end{psmallmatrix}
\in \Sp(2g,\Z)$ acts on the Siegel space and on $z$ by
\[
  \Omega' = (A\Omega+B)(C\Omega+D)^{-1}, \qquad z' = (C\Omega+D)^{-\mathsf T} z .
\]

\begin{theorem}[Transformation law]\label{thm:weil}
For $k$ even, the theta space is carried to itself:
\[
  \theta_j(z';\Omega') \;=\; j(\mathsf M,\Omega)\,\exp\!\big(\pi i k\, z^{\mathsf T}(C\Omega+D)^{-1}C\,z\big)
  \sum_l \rho_k(\mathsf M)_{jl}\,\theta_l(z;\Omega),
\]
with $\rho_k(\mathsf M)$ a $k^g\times k^g$ matrix independent of $z$ and of
$\Omega$, unitary, and satisfying $\rho_k(\mathsf M_1\mathsf M_2) = (\text{phase})\,
\rho_k(\mathsf M_1)\rho_k(\mathsf M_2)$. This is the \emph{Weil representation}.
\end{theorem}

Here $j(\mathsf M,\Omega)$ includes $\det(C\Omega+D)^{1/2}$ and a compatible
metaplectic phase; it is not independent of $\Omega$. For genus-one $S$,
choose $j(S,\tau)=(-i\tau)^{1/2}$. The fiber norm and lift convention are
part of the normalized identification. Odd levels require compatible
characteristics or a restricted modular group; the displayed even-level
generator formulas do not apply unchanged. The minimal fixed-basis odd-level
implementation uses the theta subgroup: in genus one test $S$ and $T^2$,
and in higher genus only generators preserving the characteristic (a shear
with even diagonal is allowed). Reject $T$ at odd level in that mode.
Alternatively implement and track the transformed characteristics and
line bundles explicitly. An incompatible fit is a failed hypothesis,
not numerical noise.

Two generators suffice for $g = 1$. For $T = \begin{psmallmatrix}1&1\\0&1\end{psmallmatrix}$
($\tau\mapsto\tau+1$) the sum changes by $e^{\pi i k(n+j/k)^2}$, and for $k$
even this is $e^{\pi i j^2/k}$, independent of $n$:
\[
  \rho_k(T) = \operatorname{diag}\big(e^{\pi i j^2/k}\big)_{j}, \qquad
  \rho_2(T) = \operatorname{diag}(1, i).
\]
For $S = \begin{psmallmatrix}0&-1\\1&0\end{psmallmatrix}$ ($\tau\mapsto -1/\tau$)
Poisson summation gives
\[
  \rho_k(S)_{jl} = \frac{1}{\sqrt{k}}\, e^{-2\pi i jl/k}, \qquad
  \rho_2(S) = \frac{1}{\sqrt2}\begin{pmatrix}1&1\\1&-1\end{pmatrix},
\]
the Hadamard gate. In the code $\rho_k(\mathsf M)$ is never typed in: it is fitted
from the transformation law at sampled $z$. The residual is a numerical
check, not a proof of closure for all $z$; use held-out samples and truncation
convergence. The formulas above are independent checks.

A \emph{shear} $\begin{psmallmatrix}I&B\\0&I\end{psmallmatrix}$ with $B$
symmetric integer sends $\Omega\mapsto\Omega+B$ and acts diagonally by
$e^{\pi i\, j^{\mathsf T}Bj/k}$. For $g = 2$, $B = \begin{psmallmatrix}0&1\\1&0\end{psmallmatrix}$,
$k = 2$: the phase is $(-1)^{j_1 j_2}$, the controlled-$Z$ gate.

\begin{proposition}[Products]\label{prop:kron}
$\rho_k(\mathsf M_1\oplus \mathsf M_2) = \rho_k(\mathsf M_1)\otimes\rho_k(\mathsf M_2)$. A block-diagonal
monodromy is a product operator (operator Schmidt rank one) and cannot
entangle; the shear above has Schmidt rank two.
\end{proposition}

Because $\rho_k(\mathsf M)$ does not depend on $\Omega$, the metric cannot deform
it: on this register ``the metric changes the representative, not the
operator'' holds verbatim (different $\tau$ are projectively equivalent
polarizations under the declared projectively flat identification).
This identifies coefficient spaces; it does not assert that evaluation
coherent states at arbitrary corresponding real holonomies are parallel.

\begin{proposition}[Why compatible geometric transport is unitary]\label{prop:geometric-unitarity}
A compatible scalar lift of a symplectic cycle relabelling acts on the
finite Heisenberg group preserving its center and adjoint. Its action
on the theta register is implemented by a unitary operator unique up to
phase. Compatible successive lifts compose projectively.
\end{proposition}
\begin{proof}
The clock operators have a common eigenvector; its shift orbit gives the
$d$ distinct joint eigencharacters and spans an irreducible representation.
The same construction in the relabelled representation gives an
intertwiner $U$, unique up to scalar because the translations span the
full matrix algebra. Preservation of adjoint implies $U^\dagger U$
commutes with every translation, hence is a positive scalar identity.
Normalize that scalar to obtain a unitary. Composing intertwiners gives
another intertwiner, so differs only by phase. The scalar lift and
characteristic compatibility are hypotheses, especially at odd level.
\end{proof}

\section{Orientation reversal}\label{sec:antiunitary}

An orientation-reversing period map satisfies
$M_{\rm per}^{\mathsf T}JM_{\rm per}=-J$. Put
$R_{\rm per}=\operatorname{diag}(I,-I)$ and
$N=M_{\rm per}R_{\rm per}$, so $N$ is symplectic.
First apply $R_{\rm per}$ to the period coordinates $(a,b)\mapsto(a,-b)$:
\[
 \Omega\mapsto-\bar\Omega,\qquad z\mapsto-\bar z,\qquad
 \theta_j(-\bar z;-\bar\Omega)=\overline{\theta_j(z;\Omega)}.
\]
Then apply the orientation-preserving bridge of
Proposition~\ref{prop:coordinate-bridge} to $N$. With
$\mathsf N$ its Siegel matrix the value-column action is
$\rho_k(\mathsf N)K$, where $K$ is complex conjugation.
The antiunitary lift is part of the data; the orientation-preserving
bridge is not blindly applied to an anti-symplectic matrix.
For the torus period swap, $N=S$ and $\mathsf N=S^{-1}$, giving Hadamard
followed by conjugation at level two. An antiunitary is not a CP channel;
orientation reversal alone does not derive physical time reversal.

\section{Coherent states and their exact covariance}\label{sec:coherent}

\begin{definition}[Evaluation coherent state]
In the normalized theta basis, set $F_\Omega(z)=(\theta_j(z;\Omega))_j$ and
\[
 c_\Omega(z)=\frac{\overline{F_\Omega(z)}}{\|F_\Omega(z)\|_2}.
\]
This is the Riesz representative of evaluation, up to its local line-fiber
scalar. Its ray is well defined wherever evaluation is nonzero; reject a
zero evaluation rather than divide by zero. A change of local line
trivialization multiplies the vector by a phase after normalization.
\end{definition}

\begin{proposition}[Covariance with the value-column convention]\label{prop:coherent-covariance}
If the theta transformation law is
$F_{\Omega'}(z')=f(z,\Omega)\rho_k(\mathsf M)F_\Omega(z)$ with unitary
$\rho_k(\mathsf M)$ and $f\ne0$, then
\[
 c_{\Omega'}(z')=\frac{\bar f}{|f|}
                 \overline{\rho_k(\mathsf M)}\,c_\Omega(z).
\]
For the antiunitary value action $\rho_k(\mathsf N)K$, the coherent
action is $\overline{\rho_k(\mathsf N)}K$, with the same scalar convention.
\end{proposition}
\begin{proof}
Conjugate the value-column law and divide by its norm. Unitarity removes
the matrix from the norm; only $|f|$ remains.
\end{proof}
One can instead define the state representation to be the conjugate
representation from the outset. Either convention works; fitting the
value representation and applying it unaltered to conjugated evaluations
does not. The two representations have identical algebraic completeness.

For a general flat cobordism, the period relation gives
$(a_{\rm out},b_{\rm out})=M_{\rm per}(a_{\rm in},b_{\rm in})$ modulo
integers. This alone gives no relation between $\Omega_{\rm out}$ and
$\Omega_{\rm in}$. Measure the holonomy relation separately. Assert the
displayed coherent covariance only when the polarization and coordinate
hypotheses hold. Different-modulus identity collars are negative controls:
projectively flat identification of coefficient spaces is not a theorem
that their evaluation states coincide.

\begin{theorem}[Complete single-qubit coherent encoding]\label{thm:qubit-coverage}
For any elliptic curve $E$ and positive degree-two line bundle $L^2$,
the coherent-state map $E\to\mathbb{CP}^1$ is onto. Thus every pure qubit
ray is encoded by some real holonomy pair, at fixed polarization.
\end{theorem}
\begin{proof}
Riemann--Roch on a genus-one curve gives $h^0(L^2)=2$ and
$h^0(L^2(-p))=1$ for every point $p$. Thus the two basis sections have
no common zero. Their ratio defines a nonconstant holomorphic map
$E\to\mathbb{CP}^1$; compactness and the open mapping theorem make its
image all of $\mathbb{CP}^1$. Conjugating coordinates preserves
surjectivity. The map has degree two, not a unique inverse.
\end{proof}
This proves existence of an encoding, not physical preparation or a
spontaneous choice of state. At fixed polarization a single elliptic
curve cannot cover all qutrit rays in $\mathbb{CP}^2$. More generally the
coherent family has complex dimension at most $g$, compared with
$k^g-1$ for all rays. Arbitrary state projectors remain representable by
Corollary~\ref{cor:geometric-states}; coherent evaluation is only one family.

\section{Seams, factorization and entangling capacity}\label{sec:seam}

Take two collars $W_1, W_2$ (each $T^2\times[0,1]$) and join them.

\begin{proposition}[Sphere joins cannot entangle]\label{prop:sphere}
If the join is along a $2$-sphere (remove a ball from each and identify the
boundary spheres --- a connected sum), then $H^1(W) = H^1(W_1)\oplus H^1(W_2)$,
the monodromy is block-diagonal $M_1\oplus M_2$, and $\rho_k(M) =
\rho_k(M_1)\otimes\rho_k(M_2)$ for every choice of $M_1, M_2$.
\end{proposition}
\begin{proof}
Mayer--Vietoris with $W_1\cap W_2 \simeq S^2$: $H^0(W_1)\oplus H^0(W_2)\to
H^0(S^2)$ is onto and $H^1(S^2) = 0$, so $H^1(W)\to H^1(W_1)\oplus H^1(W_2)$
is an isomorphism. The external period relation therefore splits with the
original markings, even if harmonic representatives have support in both
pieces. Proposition~\ref{prop:kron} finishes it.
\end{proof}

For valid gluing with overlap retracting to $S^1$, Mayer--Vietoris imposes
equality of restricted periods in $H^1(S^1;\C)$. This may couple classes,
but guarantees neither a full-domain graph, an invertible symplectic
monodromy nor an entangling gate. An ordinary $3$-dimensional $1$-handle
does not automatically provide this overlap; specify its attaching regions
and verify the manifold and boundary topology.
Mixing factors is not sufficient: SWAP has nontrivial operator Schmidt rank
but maps every product state to a product state. Controlled-$Z$ does
entangle, sending $|+\rangle|+\rangle$ to
$(|00\rangle+|01\rangle+|10\rangle-|11\rangle)/2$, with reduced state $I/2$.
This witness distinguishes entangling capacity from operator-vector correlations.

\begin{remark}[Measured after consolidation]\label{rem:seam-measured}
Three constructions have since been read (Section~\ref{sec:evidence-merged}).
A $1$-handle between the far tori of two collars (a tube over one triangle
of each) makes the far boundary one genus-two surface; its monodromy stays
block-diagonal (the tube alone couples no cycle) while the surface's own
period matrix acquires an off-diagonal entry $\Omega_{12}$ set by the
neck's waist and length --- the metric coupling of Section~\ref{sec:tensor},
and sampled coherent states have nonzero entropy in the declared
two-factor register. Nonzero $\Omega_{12}$ alone is not a theorem that
every evaluation point is entangled, or that the identity monodromy is
an entangling operator. Layered flip passes on a collar's far torus glue one
cobordism on another along the far surface and read the Dehn twists
$T_A$, $T_B$ and the quarter turn $S$, with passes stacked in order
composing and the two orders of a non-commuting pair giving different
matrices. The collar of the $\Z/10$-symmetric genus-two surface relabelled
by its order-five rotation reads an integer symplectic monodromy of order
five that mixes the two factors and acts on the level-two register as a
two-qubit Clifford gate sending product stabilizer states to Bell pairs:
an entangling monodromy realized geometrically, by the product-input
witness above, not merely a mixing matrix. A Dehn twist along a curve
through the tube's neck --- the controlled-$Z$ of the shear read on a
torus-pair register --- has no fixed flip pass landing on an isomorphic
far triangulation and remains open.
\end{remark}

\section{Clifford, and its limit}\label{sec:clifford}

Weil operators normalize the finite Heisenberg--Weyl group and are
Clifford operators. State explicitly which translations and central
phases are included in a chosen lifted image; do not assume that it is a
proper subgroup of the projective Clifford group. At even level lift and
congruence conventions can involve $2k$,
so naive identification with $\Sp(2g,\Z/k)$ loses phase data. At $k=2$,
these operators with stabilizer preparations and measurements are efficiently
simulable and not quantum-universal. Metric-dependent transfer is not
automatically a non-Clifford unitary. Non-abelian quantization is one
proposed resource, not an exhaustive list of routes to universality.

The finite gate-group restriction does not contradict
Theorem~\ref{thm:full-algebra}: the complex linear span of geometric
translations is the entire matrix algebra, while products of a restricted
set of unitary gates need not be dense. The experiment targets the algebra
representation, not computational universality or efficient synthesis.

\section{From a Lie algebra to its root system}\label{sec:lie}

Everything non-abelian below is generated by one piece of structure: the
eigenvalues of a maximal commuting family inside a Lie algebra. Fix a
compact simple Lie algebra $\mathfrak g$ ($\mathfrak{su}(2)$,
$\mathfrak{su}(3)$, \dots) and a maximal abelian subalgebra
$\mathfrak h\subset\mathfrak g$, the \emph{Cartan subalgebra}; its
dimension $r$ is the \emph{rank}. Complexify. The operators
$\operatorname{ad}h = [h,\cdot\,]$, $h\in\mathfrak h$, commute and are
diagonalizable, so
\[
  \mathfrak g_\C = \mathfrak h_\C\oplus\bigoplus_{\alpha\in\Phi}\mathfrak g_\alpha,
  \qquad [h,x] = \alpha(h)\,x \quad (x\in\mathfrak g_\alpha),
\]
with $\Phi\subset\mathfrak h_\C^*$ a finite set, the \emph{roots}; each
$\mathfrak g_\alpha$ is one-dimensional and $\Phi = -\Phi$. The roots are
the weights of the adjoint representation. The Killing form
$\kappa(x,y) = \Tr(\operatorname{ad}x\operatorname{ad}y)$ is nondegenerate
on $\mathfrak h$. On the real span of roots in the complexified dual use
the induced positive Euclidean form (the compact real form has negative
Killing form), normalized throughout so that long roots have
$|\alpha|^2 = 2$. The \emph{coroot} is $\alpha^\vee = 2\alpha/|\alpha|^2$.
A generic hyperplane splits $\Phi = \Phi^+\sqcup(-\Phi^+)$; the
\emph{simple roots} $\alpha_1,\dots,\alpha_r$ are the positive roots that
are not sums of two positive roots, every positive root is a non-negative
integer combination of them, and the \emph{Cartan matrix}
$A_{ij} = \langle\alpha_i,\alpha_j^\vee\rangle$ has $A_{ii} = 2$ and
$A_{ij}\in\{0,-1,-2,-3\}$ off the diagonal. The \emph{highest root}
$\theta$ is the positive root with the largest coefficients; it is long,
so $\theta^\vee = \theta$.

Three lattices in $\mathfrak h^*$:
\[
 \begin{aligned}
  Q &= \Z\Phi &&\text{(root lattice)},\\
  Q^\vee &= \Z\Phi^\vee &&\text{(coroot lattice)},\\
  P &= \{\mu : \langle\mu,\alpha^\vee\rangle\in\Z\ \ \forall\alpha\in\Phi\}
       =(Q^\vee)^* &&\text{(weight lattice)}.
 \end{aligned}
\]
$P$ is spanned by the \emph{fundamental weights} $\omega_i$,
$\langle\omega_i,\alpha_j^\vee\rangle = \delta_{ij}$, so
$\alpha_i = \sum_jA_{ij}\omega_j$ and $Q\subset P$ has index $\det A$, the
order of the centre of the simply connected group. The \emph{Weyl vector}
is $\rho = \tfrac12\sum_{\alpha>0}\alpha = \sum_i\omega_i$ and the
\emph{dual Coxeter number} is $h^\vee = \langle\rho,\theta^\vee\rangle+1$.

\begin{example}[$\mathfrak{su}(2)$; Problem P28]\label{ex:sl2}
$\mathfrak g_\C = \mathfrak{sl}(2,\C)$ with $h = \operatorname{diag}(1,-1)$,
$e = E_{12}$, $f = E_{21}$: $[h,e] = 2e$, $[h,f] = -2f$, $[e,f] = h$. One
positive root $\alpha$, $\alpha(h) = 2$, $\mathfrak g_\alpha = \C e$;
$\alpha^\vee = h$; the fundamental weight has $\omega(h) = 1$, so
$\alpha = 2\omega$, $P = \Z\omega$, $Q = 2\Z\omega$, $P/Q = \Z/2$,
$|\omega|^2 = \tfrac12$, $\rho = \omega$, $\theta = \alpha$, $h^\vee = 2$,
$A = (2)$, $W = \{\pm1\}$.
\end{example}

\begin{example}[$\mathfrak{su}(3)$; Problem P29]\label{ex:sl3}
$\mathfrak{sl}(3,\C)$ with $\mathfrak h$ the traceless diagonal matrices
and $\varepsilon_i(\operatorname{diag}(a_1,a_2,a_3)) = a_i$. The roots are
$\varepsilon_i-\varepsilon_j$, $i\ne j$, with root vectors $E_{ij}$: six
roots. Simple roots $\alpha_1 = \varepsilon_1-\varepsilon_2$,
$\alpha_2 = \varepsilon_2-\varepsilon_3$; the third positive root
$\theta = \alpha_1+\alpha_2$ is the highest. The Cartan matrix is
$\begin{psmallmatrix}2&-1\\-1&2\end{psmallmatrix}$: the simple roots have
length $\sqrt2$ at $120^\circ$, and the six roots are the vertices of a
regular hexagon --- the root system $A_2$, whose root lattice is the
hexagonal lattice. $\omega_1 = \tfrac13(2\alpha_1+\alpha_2)$,
$\omega_2 = \tfrac13(\alpha_1+2\alpha_2)$, $P/Q = \Z/3$ (triality),
$\rho = \omega_1+\omega_2 = \theta$, $h^\vee = 3$. The Weyl group permutes
the $\varepsilon_i$: $W = S_3\cong D_3$.
\end{example}

\begin{proposition}[An internal $A_2$ lattice from scalar lengths]\label{prop:scalar-a2}
Let a Euclidean lattice $Q=\Z\alpha_1+\Z\alpha_2$ have a generating
triangle with side lengths $\sqrt2,\sqrt2,\sqrt6$ at
$0,\alpha_1,\alpha_2$. Its Gram matrix, recovered by the cosine law, is
\[
 A=\begin{pmatrix}2&-1\\-1&2\end{pmatrix}.
\]
The shortest nonzero lattice vectors are
$\pm\alpha_1,\pm\alpha_2,\pm(\alpha_1+\alpha_2)$; they form $A_2$.
Their reflections preserve $Q$, generate $S_3$, and the metric-dual
lattice is $P=A^{-1}\Z^2$ in root coordinates, with $|P/Q|=3$.
\end{proposition}
\begin{proof}
$\alpha_1\cdot\alpha_2=(2+2-6)/2=-1$. The squared norm of
$m\alpha_1+n\alpha_2$ is $2(m^2-mn+n^2)$, whose smallest positive
value is two at exactly the six stated integer pairs. Reflection in
$\alpha_i$ is $v\mapsto v-(v\cdot\alpha_i)\alpha_i$; the integral
Gram entries prove lattice preservation, and the two reflections have
product of order three. Duality gives $P$ and $\det A=3$.
\end{proof}
This internal lattice is chosen independently of the spacetime boundary.
It proves a scalar-geometric representation of the root datum, not that
a generic boundary triangulation spontaneously has this symmetry, nor
that its harmonic mass Gram automatically equals $A$.

\begin{theorem}[Highest-weight classification]\label{thm:highest}
The finite-dimensional irreducible representations of $\mathfrak g$ are
$V_\lambda$, one for each dominant weight
$\lambda\in P^+ = \{\lambda\in P : \langle\lambda,\alpha_i^\vee\rangle\ge0\ \forall i\}$.
$V_\lambda = \bigoplus_{\mu\in P}V_\lambda[\mu]$ decomposes into weight
spaces, $h$ acting on $V_\lambda[\mu]$ by $\mu(h)$; $\lambda$ occurs once,
every other weight is $\lambda$ minus a non-negative combination of simple
roots, and the weights counted with multiplicity form a $W$-invariant
multiset.
\end{theorem}

For $\mathfrak{su}(2)$, $V_\lambda = \operatorname{Sym}^\lambda\C^2$ (spin
$\lambda/2$) has weights $\lambda,\lambda-2,\dots,-\lambda$ in units of
$\omega$, each once: dimension $\lambda+1$. The \emph{character}
$\operatorname{ch}V_\lambda = \sum_\mu\dim V_\lambda[\mu]\,e^{\mu}$ is a
function on the maximal torus $T = \exp\mathfrak h$ through
$e^{\mu}(\exp h) = e^{\mu(h)}$; it is exactly the object of
Section~\ref{sec:weights}: characters multiply under $\otimes$ because
weights add.

\section{The Weyl group and the character formula}\label{sec:weyl}

The reflection in the root $\alpha$ is
$s_\alpha(\mu) = \mu-\langle\mu,\alpha^\vee\rangle\alpha$. The \emph{Weyl
group} $W = \langle s_\alpha:\alpha\in\Phi\rangle$ is finite, generated by
the simple reflections $s_i = s_{\alpha_i}$, permutes $\Phi$, preserves
$P$, $Q$, $Q^\vee$ and $\langle\cdot,\cdot\rangle$, and has the closed
dominant chamber $\{\langle\mu,\alpha_i^\vee\rangle\ge0\}$ as a
fundamental domain. Its sign character is
$\varepsilon(w) = \det w = (-1)^{\ell(w)}$.

\begin{theorem}[Weyl character formula; Problem P31]\label{thm:weylchar}
For $\lambda\in P^+$,
\[
  \operatorname{ch}V_\lambda = \frac{\sum_{w\in W}\varepsilon(w)\,e^{w(\lambda+\rho)}}
                                    {\sum_{w\in W}\varepsilon(w)\,e^{w\rho}},
  \qquad
  \sum_{w\in W}\varepsilon(w)\,e^{w\rho} = e^{\rho}\prod_{\alpha>0}\big(1-e^{-\alpha}\big),
\]
and $\dim V_\lambda = \prod_{\alpha>0}\langle\lambda+\rho,\alpha^\vee\rangle/\langle\rho,\alpha^\vee\rangle$.
\end{theorem}

For $\mathfrak{su}(2)$, writing $x = e^{\omega}$:
$\operatorname{ch}V_\lambda = (x^{\lambda+1}-x^{-\lambda-1})/(x-x^{-1})
= x^{\lambda}+x^{\lambda-2}+\dots+x^{-\lambda}$. The numerator is the
\emph{antisymmetrization} of the single exponential $e^{\lambda+\rho}$
under $W$; every formula below has that shape.

\begin{proposition}[Brauer--Klimyk; weights add, the Weyl group folds; Problem P32]
\label{prop:klimyk}
Let $m_\mu(\nu) = \dim V_\mu[\nu]$. Then, as virtual representations,
\[
  V_\lambda\otimes V_\mu \;=\; \sum_{\nu\in P}m_\mu(\nu)\;\varepsilon(w_\nu)\;V_{w_\nu(\lambda+\nu+\rho)-\rho},
\]
where $w_\nu\in W$ takes $\lambda+\nu+\rho$ into the dominant chamber and
the term is omitted when $\lambda+\nu+\rho$ lies on a wall of the chamber.
The right side has non-negative coefficients.
\end{proposition}

\begin{proof}
Multiply the numerator of $\operatorname{ch}V_\lambda$ by
$\operatorname{ch}V_\mu = \sum_\nu m_\mu(\nu)e^{\nu}$. Since $m_\mu$ is
$W$-invariant,
$\sum_w\varepsilon(w)e^{w(\lambda+\rho)}\sum_\nu m_\mu(\nu)e^{w\nu}
= \sum_\nu m_\mu(\nu)\sum_w\varepsilon(w)e^{w(\lambda+\nu+\rho)}$, and each
inner sum is $\varepsilon(w_\nu)$ times the numerator of
$\operatorname{ch}V_{w_\nu(\lambda+\nu+\rho)-\rho}$, or zero when
$\lambda+\nu+\rho$ is fixed by a reflection. Divide by the denominator.
\end{proof}

For $\mathfrak{su}(2)$: $V_1\otimes V_\lambda = V_{\lambda+1}\oplus V_{\lambda-1}$
for $\lambda\ge1$ (Clebsch--Gordan), and $V_1\otimes V_0 = V_1$ because
$\lambda+\nu+\rho = 0$ for $\nu = -1$ lies on the wall. This is the finite
prototype of the Kac--Walton rule of Section~\ref{sec:verlinde}: at level
$k$ one more wall appears and one more reflection folds across it.

\section{Affine Lie algebras: the level, the translations, and theta functions}\label{sec:affine}

The register needs the loop version of $\mathfrak g$. The \emph{affine Lie
algebra} is
\[
  \hat{\mathfrak g} = \mathfrak g[t,t^{-1}]\oplus\C K\oplus\C d,\qquad
  [x\otimes t^m,\,y\otimes t^n] = [x,y]\otimes t^{m+n} + m\,\delta_{m+n,0}\langle x,y\rangle K,
\]
with $K$ central and $d = t\,\tfrac{d}{dt}$ the grading. A representation
on which $K$ acts by the scalar $k$ has \emph{level} $k$. The Cartan
subalgebra grows to $\hat{\mathfrak h} = \mathfrak h\oplus\C K\oplus\C d$,
a weight is $\hat\lambda = (\lambda,k,n)$, and there is one new simple root
$\alpha_0 = \delta-\theta$, where $\delta$ generates the imaginary roots
($\delta(d) = 1$, $\delta|_{\mathfrak h\oplus\C K} = 0$). Dominance for
$\alpha_0$ reads $\langle\hat\lambda,\alpha_0^\vee\rangle = k-\langle\lambda,\theta^\vee\rangle\ge0$:

\begin{theorem}[Kac]\label{thm:integrable}
The integrable highest-weight modules $L(\hat\lambda)$ of level $k$ are
indexed by the \emph{Weyl alcove}
$P_k^+ = \{\lambda\in P^+ : \langle\lambda,\theta^\vee\rangle\le k\}$.
\end{theorem}

The \emph{affine Weyl group} is $\hat W = W\ltimes t(Q^\vee)$: the finite
Weyl group together with translations $t_\gamma$, $\gamma\in Q^\vee$, which
act on a weight of level $K$ by
\[
  t_\gamma(\hat\lambda) = \hat\lambda + K\gamma - \big(\langle\lambda,\gamma\rangle+\tfrac K2|\gamma|^2\big)\delta,
\]
i.e.\ by $\lambda\mapsto\lambda+K\gamma$ on the finite part. It is
generated by $s_1,\dots,s_r$ and the affine reflection
$s_0 = t_{\theta^\vee}s_\theta$, the reflection in the hyperplane
$\langle\mu,\theta^\vee\rangle = K$ (Problem P34). At level $K$ the
fundamental alcove
$\{\mu : \langle\mu,\alpha_i^\vee\rangle\ge0,\ \langle\mu,\theta^\vee\rangle\le K\}$
is a simplex and a fundamental domain, and $P_k^+ + \rho$ is its set of
interior lattice points at $K = k+h^\vee$.

\begin{theorem}[Weyl--Kac; the lattice sum is the translation orbit; Problem P35]\label{thm:weylkac}
With $\hat\rho = (\rho,h^\vee,0)$,
\[
  \operatorname{ch}L(\hat\lambda) = \frac{\sum_{\hat w\in\hat W}\varepsilon(\hat w)\,e^{\hat w(\hat\lambda+\hat\rho)-\hat\rho}}
  {\prod_{\hat\alpha>0}(1-e^{-\hat\alpha})^{\operatorname{mult}\hat\alpha}} .
\]
Write $q = e^{-\delta} = e^{2\pi i\tau}$ and $e^{\mu} = e^{2\pi i\langle\mu,z\rangle}$,
and split $\hat W = W\ltimes t(Q^\vee)$. The translation orbit of a single
exponential at level $K$ is a theta function,
\[
  \sum_{\gamma\in Q^\vee}e^{t_\gamma(\hat\mu)} = q^{-|\mu|^2/2K}\,\Theta^{(K)}_\mu(\tau,z),\qquad
  \Theta^{(K)}_\mu(\tau,z) = \sum_{\gamma\in Q^\vee}q^{\frac K2|\frac\mu K+\gamma|^2}\,e^{2\pi iK\langle\frac\mu K+\gamma,\,z\rangle},
\]
so the numerator at $K = k+h^\vee$ is
$\sum_{w\in W}\varepsilon(w)\Theta^{(K)}_{w(\lambda+\rho)}$ up to a power of
$q$; the denominator is the same expression at $K = h^\vee$, $\mu = \rho$
(the affine denominator identity); and the normalized character is the
ratio of Theorem~\ref{thm:kac}.
\end{theorem}

This is the structural content of ``geometry is algebra'' in the
non-abelian register. The translations $t(Q^\vee)$ are a lattice, and the
lattice sum they generate is a theta function --- the object
Section~\ref{sec:theta} already attaches to a torus. The finite Weyl
group $W$ antisymmetrizes the sum, but does not replace the root/coroot
lattices, their normalization or the level. An abelian theta register has
trivial Weyl group and no Weyl denominator; substituting $h^\vee=0$ into
the non-abelian character ratio would leave its denominator undefined.

\begin{theorem}[Sugawara; Problem P36]\label{thm:sugawara}
On a level-$k$ integrable module the Virasoro generator
$L_0 = \tfrac{1}{2(k+h^\vee)}\sum_a{:}J^aJ^a{:}$ acts on the highest-weight
vector of $L(\hat\lambda)$ by the \emph{conformal weight}
$h_\lambda = C_2(\lambda)/2(k+h^\vee)$,
$C_2(\lambda) = \langle\lambda,\lambda+2\rho\rangle = |\lambda+\rho|^2-|\rho|^2$,
with \emph{central charge} $c = k\dim\mathfrak g/(k+h^\vee)$. The
$T$-matrix of the modular data is $T_\lambda = e^{2\pi i(h_\lambda-c/24)}$.
\end{theorem}

With the strange formula $|\rho|^2/2h^\vee = \dim\mathfrak g/24$ this is
$T_\lambda = \exp2\pi i\big[|\lambda+\rho|^2/2(k+h^\vee)-\dim\mathfrak g/24\big]$,
the form quoted in Theorem~\ref{thm:kp}. For $\mathfrak{su}(2)_k$:
$h_\lambda = \lambda(\lambda+2)/4(k+2)$ (spin $j = \lambda/2$:
$j(j+1)/(k+2)$) and $c = 3k/(k+2)$; $c = 1$ at $k = 1$, $c = \tfrac32$
at $k = 2$. This has Ising-type fusion, but is not the standard Ising
conformal theory, whose central charge is $1/2$; twists also distinguish them.

\section{Root systems of rank two}\label{sec:roots}

A boundary torus supplies a rank-two integral lattice with an alternating
intersection form, not a positive root Gram or a Lie algebra. If a rank-two
Euclidean root system is additionally chosen, the possibilities are
$A_1\times A_1$, $A_2$, $B_2$, $G_2$. The complex dimension of the base
elliptic curve differs from the rank of the Cartan fiber; affine characters
on a torus can have a Lie algebra of any rank. With long roots of length $\sqrt2$
and $\alpha_1$ long (Problem P30):
\[
\begin{array}{c|cccc}
 & A_1\times A_1 & A_2 & B_2 & G_2\\ \hline
\text{Cartan } A & \begin{psmallmatrix}2&0\\0&2\end{psmallmatrix} & \begin{psmallmatrix}2&-1\\-1&2\end{psmallmatrix}
  & \begin{psmallmatrix}2&-2\\-1&2\end{psmallmatrix} & \begin{psmallmatrix}2&-3\\-1&2\end{psmallmatrix}\\[8pt]
\text{Gram }\langle\alpha_i,\alpha_j\rangle & \begin{psmallmatrix}2&0\\0&2\end{psmallmatrix} & \begin{psmallmatrix}2&-1\\-1&2\end{psmallmatrix}
  & \begin{psmallmatrix}2&-1\\-1&1\end{psmallmatrix} & \begin{psmallmatrix}2&-1\\-1&\tfrac23\end{psmallmatrix}\\[8pt]
\text{root lattice} & \text{square} & \text{hexagonal} & \text{square} & \text{hexagonal}\\
W & D_2\ (4) & D_3\ (6) & D_4\ (8) & D_6\ (12)\\
P/Q & \Z/2\times\Z/2 & \Z/3 & \Z/2 & 1\\
h^\vee & (2,2) & 3 & 3 & 4
\end{array}
\]
$W(A_1\times A_1)\subset W(B_2)$ and $W(A_2)\subset W(G_2)$ have index two;
For $A_1\times A_1$, levels and dual Coxeter numbers are specified per
simple factor. The abelian case
$\mathfrak g = \mathfrak u(1)$ has no roots, trivial $W$, $P = \Z$: the
theta functions of Section~\ref{sec:theta} are the $r = 1$, $W = 1$
abelian construction, with no non-abelian denominator.

\begin{proposition}[Rank-two Weyl groups are the symmetries a torus can have; Problem P38]
\label{prop:pointgroups}
The finite subgroups of $\GL(2,\Z)$ are the crystallographic point groups:
cyclic of order $1,2,3,4,6$ and dihedral $D_1,D_2,D_3,D_4,D_6$. The dihedral
ones $D_2,D_3,D_4,D_6$ are the four rank-two Weyl groups. A finite group of
simplicial automorphisms of a torus induces such a subgroup on $H_1$.
Translations can lie in the kernel, so the full simplicial automorphism
group need not embed in $\GL(2,\Z)$.
\end{proposition}

The square grid with one diagonal realizes only $\{\pm I,\pm\sigma\}\cong D_2$
--- the diagonal breaks the fourfold symmetry --- which is abelian, so
composition order is invisible there. The triangular (hexagonal) lattice
realizes $D_6 = W(G_2)$, non-abelian, containing $D_3 = W(A_2)$.

\section{The alcove and the Verlinde space}\label{sec:alcove}

At level $k$ the analogue of the characteristics $j\in\Z/k$ is the
\emph{Weyl alcove} of Theorem~\ref{thm:integrable},
\[
  P_k^+ = \{\lambda\in P^+ : \langle\lambda,\theta^\vee\rangle \le k\},
\]
the dominant weights of level at most $k$. For $\mathfrak{su}(2)$:
$\lambda = 0,1,\dots,k$ (spin $0,\tfrac12,\dots,\tfrac k2$), so $k+1$ states.
The torus Hilbert space is spanned by one function per alcove weight: the
\emph{Weyl--Kac character} of the level-$k$ integrable representation of the
affine algebra $\hat{\mathfrak g}$.

\begin{theorem}[Weyl--Kac; characters are antisymmetrized theta functions]
\label{thm:kac}
For $\mu\in P$ and a level $K$, set
\[
  \Theta^{(K)}_\mu(\tau,z) = \sum_{\gamma\in Q^\vee}
  \exp\!\Big(\pi i K\big|\tfrac{\mu}{K}+\gamma\big|^2\tau
  + 2\pi i K\big\langle\tfrac{\mu}{K}+\gamma, z\big\rangle\Big),
  \qquad
  A^{(K)}_\mu = \sum_{w\in W}\varepsilon(w)\,\Theta^{(K)}_{w\mu}.
\]
Then the level-$k$ character with highest weight $\lambda\in P_k^+$ is
\[
  \chi_\lambda(\tau,z) \;=\; \frac{A^{(k+h^\vee)}_{\lambda+\rho}(\tau,z)}
                                  {A^{(h^\vee)}_{\rho}(\tau,z)} .
\]
\end{theorem}

This is Theorem~\ref{thm:weylkac} written out: the sum over $Q^\vee$ is the
translation part of $\hat W$, the sum over $W$ its finite part (Problems
P41, P42). For $\mathfrak{su}(2)$ ($r=1$, $W=\{\pm1\}$, $h^\vee = 2$) this is fully
explicit in the theta functions of Definition~\ref{def:theta} (with the
variable $z$ rescaled by the root normalization):
\[
  \chi_\lambda(\tau,z) = \frac{\theta^{(2(k+2))}_{\lambda+1}(z;\tau)
      - \theta^{(2(k+2))}_{-(\lambda+1)}(z;\tau)}
     {\theta^{(4)}_{1}(z;\tau) - \theta^{(4)}_{-1}(z;\tau)},
  \qquad \lambda = 0,\dots,k,
\]
where $\theta^{(K)}_m$ is the level-$K$ theta function of characteristic
$m \bmod K$. The register's construction is: root data $\to$ alcove $\to$
antisymmetrize the existing lattice sum under $W$ $\to$ divide by the
denominator. A disjoint union of tori is still a tensor product, with basis
$(\lambda_1,\lambda_2)$.

\subsection{Coordinates and polarization of the root-lattice register}\label{sec:root-coordinates}
For the simply-laced controls $A_1,A_2$, let $A$ be the coroot Gram matrix,
write the Cartan variable as $\xi$ in root coordinates, and let
$m\in\Z^r$ be the Dynkin coordinates of $\mu$. Then
\[
 \mu=A^{-1}m,\qquad c=\frac{A^{-1}m}{K},\qquad
 \Omega=\tau A,\qquad z_{\theta}=A\xi .
\]
The generalized theta sum is
\[
 \Theta_\mu^{(K)}(\tau,\xi)=
 \sum_{n\in\Z^r}
 e^{\pi iK(n+c)^{\mathsf T}\tau A(n+c)+2\pi iK(n+c)^{\mathsf T}A\xi}.
\]
Characteristics are indexed by $P/KQ^\vee$, not arbitrary unrelated
rational points. The physical Cartan torus in these coordinates has
$\xi\sim\xi+n+\tau n'$ with polarization $KA$. Consequently its
section count is $\det(A)K^r$: $2K$ for $A_1$, $3K^2$ for $A_2$.
These are not the $K^r$ sections of the principal genus-$r$ bundle.
Using $\Omega=\tau A$ is a convenient series representation, not a
license to replace the period lattice or its polarization.

Under $\tau\mapsto-1/\tau$, Poisson summation over $Q^\vee$ mixes its
dual-lattice cosets. This is not the ordinary genus-$r$ Siegel inversion:
$(\tau A)\mapsto-(\tau A)^{-1}$ gives $-A^{-1}/\tau$, not $-A/\tau$.
Under $\tau\mapsto\tau+1$, $\Omega$ changes by $A$, not by $I$.
At $S$, the numerator at level $K=k+h^\vee$ and denominator at level
$h^\vee$ have common determinant factors and their exponential quotient is
$\exp(\pi i k\,\xi^{\mathsf T}A\xi/\tau)$. Derive and fit this
root-lattice law directly. Evenness of $A_1,A_2$ permits integer $K$;
the odd-level restriction of the principal rank-one register is a
different issue.

For $A_2$ the surviving character dimension is
$|P_k^+|=(k+1)(k+2)/2$, hence three at $k=1$ and six at $k=2$.
The ratios are normalized affine characters
$\operatorname{Tr}_{L_\lambda}(q^{L_0-c/24}e^{2\pi i(\cdot,\xi)})$.
On a Weyl wall numerator and denominator may both vanish: use a
controlled removable limit or evaluate away from the walls with a
denominator bound. Do not divide numerical zeros.

\section{Kac--Peterson, Verlinde, and multiplication into addition}\label{sec:verlinde}

\begin{theorem}[Kac--Peterson]\label{thm:kp}
The characters transform under $\SL(2,\Z)$ among themselves,
$\chi_\lambda(-1/\tau, z/\tau) = (\cdots)\sum_\mu S_{\lambda\mu}\chi_\mu(\tau,z)$
and $\chi_\lambda(\tau+1,z) = T_{\lambda\lambda}\chi_\lambda(\tau,z)$, with
\[
  S_{\lambda\mu} = c\sum_{w\in W}\varepsilon(w)\,
     \exp\!\Big(-\tfrac{2\pi i}{k+h^\vee}\langle w(\lambda+\rho),\mu+\rho\rangle\Big),
  \qquad
  T_{\lambda\lambda} = \exp\!\Big(2\pi i\big[\tfrac{|\lambda+\rho|^2}{2(k+h^\vee)} - \tfrac{\dim\mathfrak g}{24}\big]\Big),
\]
with $c = i^{|\Phi^+|}\,|P/(k+h^\vee)Q^\vee|^{-1/2}$ and $T$ the Sugawara
matrix of Theorem~\ref{thm:sugawara} (Problems P43, P36).
For $\mathfrak{su}(2)_k$: $S_{\lambda\mu} = \sqrt{\tfrac{2}{k+2}}\,
\sin\!\big(\tfrac{\pi(\lambda+1)(\mu+1)}{k+2}\big)$ and
$T_\lambda = \exp\!\big(2\pi i\big[\tfrac{(\lambda+1)^2}{4(k+2)} - \tfrac18\big]\big)$.
\end{theorem}

These replace the abelian Weil matrices. A numerical character experiment
can fit their transformation laws with held-out residuals, controlling
lattice truncation and removable denominator zeros. Fitting geometric lattice sums against independent character formulas
validates the representation objective. It is not evidence that a native
simplex engine generated the root data or a single fusion cobordism.

\begin{theorem}[Verlinde]\label{thm:verlinde}
The fusion coefficients $N_{\lambda\mu}^{\ \nu}$ --- the multiplicities in the
level-$k$ truncated tensor product --- are
\[
  N_{\lambda\mu}^{\ \nu} = \sum_{\sigma\in P_k^+}
  \frac{S_{\lambda\sigma}S_{\mu\sigma}\overline{S_{\nu\sigma}}}{S_{0\sigma}},
\]
and are non-negative integers. For $\mathfrak{su}(2)_k$,
$N_{\lambda\mu}^{\ \nu} = 1$ exactly when $|\lambda-\mu|\le\nu\le
\min(\lambda+\mu,\,2k-\lambda-\mu)$ and $\lambda+\mu+\nu$ is even, else $0$.
\end{theorem}

The Kac--Walton formula computes the same numbers by \emph{adding} the
weights of $V_\mu$ to $\lambda$ and reflecting every result back into the
alcove by the affine Weyl group, with the sign of the reflection. That is
the precise sense of ``root systems turn multiplication into addition'':
characters multiply, weights add, and the Weyl group folds. It is
Proposition~\ref{prop:klimyk} with $\hat W$ in place of $W$: the one new
ingredient is the affine reflection $s_0$ (Problem P44).

\begin{definition}[Fusion algebra]
Let $\mathcal V_k$ have basis $f_\lambda$, $\lambda\in P_k^+$, and product
$f_\lambda f_\mu=\sum_\nu N_{\lambda\mu}^{\ \nu}f_\nu$.
This fixed-level fusion product is not ordinary pointwise multiplication
of affine characters at a common arbitrary $(\tau,\xi)$; such a product
does not preserve the level.
\end{definition}

\begin{theorem}[Fusion as multiplication on geometric alcove points]\label{thm:fusion-evaluation}
For the unitary modular data above, put
\[
 v_\lambda(\sigma)=\frac{S_{\lambda\sigma}}{S_{0\sigma}},
 \qquad \sigma\in P_k^+.
\]
Then $f_\lambda\mapsto v_\lambda$ is an algebra isomorphism
$\mathcal V_k\otimes\C\cong\C^{P_k^+}$ with pointwise multiplication.
Moreover the Weyl formula identifies
\[
 v_\lambda(\sigma)=
 \operatorname{ch}V_\lambda
 \left(\exp\left[-\frac{2\pi i(\sigma+\rho)}{k+h^\vee}\right]\right).
\]
Thus fusion can be calculated by evaluating finite characters at specified
points of the internal geometric torus, multiplying values, and changing
back to the alcove basis.
\end{theorem}
\begin{proof}
The last identity is the ratio of the two finite Weyl alternating sums
in $S_{\lambda\sigma}$ and $S_{0\sigma}$. The evaluation matrix is
invertible because $S$ is unitary and $S_{0\sigma}>0$. Substituting the
Verlinde formula and summing over $\nu$ using unitarity gives
$\sum_\nu N_{\lambda\mu}^{\ \nu}v_\nu(\sigma)
=v_\lambda(\sigma)v_\mu(\sigma)$. The unit is $v_0=1$.
\end{proof}
The representation is exact; independent Kac--Walton folding checks the
integral nonnegative structure constants. This does not identify ordinary
representation multiplicities with affine fusion, nor prove that a
physical joining operation performs fusion.

\section{Simply-laced level one and the abelian lattice register}\label{sec:frenkelkac}

\begin{theorem}[Frenkel--Kac--Segal; Problem P37]\label{thm:fk}
Let $\mathfrak g$ be simply laced ($A$, $D$, $E$: all roots of length
$\sqrt2$, $Q = Q^\vee$). The level-one modules $L(\hat\omega_i)$, with
$\omega_i$ ranging over $0$ and the minuscule fundamental weights, are
realized on the Fock space of $r$ free bosons tensored with
$\C[Q+\omega_i]$, and their characters are the theta functions of the
cosets of the root lattice in the weight lattice:
\[
  \chi_{\omega_i}(\tau,z) = \frac{\Theta_{Q+\omega_i}(\tau,z)}{\eta(\tau)^r},\qquad
  \Theta_{Q+\omega}(\tau,z) = \sum_{\gamma\in Q+\omega}q^{|\gamma|^2/2}\,e^{2\pi i\langle\gamma,z\rangle},
\]
where $\eta(\tau) = q^{1/24}\prod_{n\ge1}(1-q^n)$.
\end{theorem}

Since $\eta$ transforms by a multiplier under $\SL(2,\Z)$, the modular
data of $\mathfrak g$ at level one is that of the \emph{even lattice} $Q$
with characteristics $P/Q$:
$S_{ab} = |P/Q|^{-1/2}e^{-2\pi i\langle a,b\rangle}$ and
$T_a = e^{-2\pi ir/24}e^{\pi i|a|^2}$ for $a,b\in P/Q$ (Poisson summation
over $Q$, whose dual is $P$). For $\mathfrak{su}(N)_1$, $P/Q = \Z/N$,
$S_{ab} = N^{-1/2}e^{2\pi iab/N}$ and $T_a\propto e^{\pi ia(N-a)/N}$. For
$\mathfrak{su}(2)_1$: $Q = \sqrt2\,\Z$, the two coset theta functions are
the level-two functions $\theta^{(2)}_0,\theta^{(2)}_1$ of
Definition~\ref{def:theta}, $S$ is the Hadamard matrix and
$T\propto\operatorname{diag}(1,i)$: the abelian register at $k = 2$ is
$\mathfrak{su}(2)$ at level one. A diagonal $T$ exists here, and not for
the rank-one register at odd level (Section~\ref{sec:weil}), because $Q$
is even: $|\gamma|^2/2\in\Z$ on $Q$, so $\tau\mapsto\tau+1$ multiplies each
coset sum by the constant $e^{\pi i|\omega_i|^2}$. The rank-one register
at level $k$ is the lattice $\sqrt k\,\Z$, even exactly when $k$ is even.
In particular $\mathfrak{su}(3)_1$ has relative twists
$(1,e^{2\pi i/3},e^{2\pi i/3})$ from the $A_2$ discriminant quadratic
form. Matching the group $\Z/3$ alone, or using the principal rank-one
level-three register, is not the Frenkel--Kac comparison.

\section{Finite on the torus, dense with punctures}\label{sec:density}

Two facts fix what a non-abelian register can and cannot do.

\begin{theorem}[Congruence property; Ng--Schauenburg]
For a modular tensor category the projective modular representation on the
unpunctured torus has a congruence kernel and finite projective image.
Exact level and linear-lift statements require fixing the normalization
of $T$ and its anomaly phases.
\end{theorem}

So on a single torus the non-abelian register, like the Clifford one, gives
a finite set of operators.

\begin{theorem}[Density in specified Jones sectors; Freedman--Larsen--Wang]
Certain Jones braid representations have projectively dense image in the
corresponding unitary group. Level, strand labels, total fusion charge and
irreducible sector must satisfy the density theorem's hypotheses.
This is not a statement for every sector of every punctured sphere.
\end{theorem}

Punctures and braids are the proposed route here, not a necessary condition
for every form of universality. Removing a bulk top cell does not by itself
construct a labeled surface puncture or conformal-block Hilbert space.
Weyl symmetry gives invariant isotypic sectors only when the operator
commutes with the action; these are not automatically level-dependent
alcove labels. The missing scalar-geometric realization and its acceptance
criteria appear in Section~\ref{sec:root-program}.

\clearpage
\part{Experiment specification}

\section{Whole harmonic states and identifiable operators}\label{sec:kernel-contract}

This section retains the native harmonic-recovery contract for existing
code and future synthesis. Its three verdicts are identifiable operator,
positive norm preservation, and requested native gate realization.
They are not prerequisites for the current theta-algebra representation
verdict in Section~\ref{sec:validation}.
A harmonic space is a candidate state space; a nonzero vector or ray needs
input data or a geometric selection rule. For positive Hermitian masses,
\[
 \langle h,L_ph\rangle_{M_p}
 =\|\delta_ph\|_{M_{p-1}}^2+\|d_ph\|_{M_{p+1}}^2.
\]
This proves closedness and coclosedness of the kernel. Orthogonality of
exact, harmonic and coexact summands gives a unique representative per
class. Complex bilinear masses do not provide this positivity proof.
For a flat torus with $\tau=a+ib$, $b>0$,
\[
 g_\tau=\frac1b\begin{pmatrix}1&a\\a&a^2+b^2\end{pmatrix},\qquad
 \omega_\tau=dx+\tau\,dy,\qquad *\omega_\tau=-i\omega_\tau .
\]
Periods give $\psi=(1,\tau)/\sqrt{1+|\tau|^2}$ under a declared identity
metric on period coordinates, not assumed Whitney orthonormality.
This fixed-orientation positive-real construction covers one Bloch
hemisphere; complex continuation needs separate diagnostics.

\subsection{Frozen whole-kernel recovery}
Assemble $K(\lambda)=\widetilde A-\lambda M_{\rm whole}$ and let $Z$ span
its geometric zero kernel in the coordinates on which readouts act.
Set $A=R_{\rm in}Z$, $B=R_{\rm out}Z$.
\begin{proposition}[Identifiability]
A unique operator $T$ on every input with $B=TA$ exists exactly when
$A$ is onto and $\ker A\subseteq\ker B$.
\end{proposition}
\begin{proof}
Necessity is immediate. Choose a right inverse $C$ of $A$. Since
$c-CAc\in\ker A$, $Bc=BCAc$; hence $T=BC$ works. Right inverses differ
by a map into $\ker A$, so give the same $T$. Onto-ness gives uniqueness.
The equivalent uniqueness test is $B(I-CA)=0$.
\end{proof}
Square invertible $A$ gives $T=BA^{-1}$. An unchecked pseudoinverse is
not a certificate. Invisible modes are allowed if both readouts annihilate
them. The zero map is valid but has no normalized Choi ray.
Record thresholds, ranks, singular gaps, conditioning and full-space
residuals. For basis and fixed-seed held-out inputs $x$, reconstruct $h$
and independently check
\[
 R_{\rm in}h=x,\qquad \widetilde A h=0,\qquad R_{\rm out}h=Tx.
\]
A minimum Euclidean image-norm witness resolves invisible modes
computationally, not physically. Frozen recovery accepts only geometry
and declared readouts, never targets or desired outputs. Rebasing $Z$
leaves $T$ unchanged; port rebasing gives
$T'=g_{\rm out}^{-1}Tg_{\rm in}$ with transformed port metrics.

\subsection{Chain/cochain and zero-mode certificates}
Distinguish the native chain frame $\Phi$ from cochain images $Z=G^U\Phi$.
For twisted boundary operators $\partial_j^U$, test
\[
 (\partial_2^{U^{-1}})^{\mathsf T}Z=0,\qquad \partial_1^U\Phi=0 .
\]
These are closedness on images and coclosedness represented on the chain
frame. Compare an independent null space with the band's rank/span;
require finite projector idempotency, rank tolerance, singular-gap and
resolvent certificates. Distinguish contour rank, algebraic multiplicity,
geometric kernel and closed/coclosed dimension on defective pencils.
Check the map to cohomology, not only equality with the Betti number.
A mismatch is a failed hypothesis, not permission to repair geometry.
Flat scalar local systems support twisted cohomology, pairing $U$ with
$U^{-1}$. Nontrivial flat $U(1)$ holonomy need not be pure gauge;
forbidding matrix connections does not remove it. Curved $U$ can have
$d_U^2\ne0$, so cohomological gluing needs new justification.

\section{Pairing and Gram conventions}\label{sec:gram-contract}

Reuse the live marked block's native dual/primal frames $F^\vee,F$,
normalized on its own Whitney mass, then embed the dual with $E$:
\[
 (F^\vee)^{\mathsf T}M_{\rm own}F=I,\qquad
 R_{\rm Gram}\Phi_W=(EF^\vee)^{\mathsf T}M_{\rm whole}\Phi_W .
\]
Reuse \texttt{PencilSchur.gramBlock}; do not duplicate mass assembly or
dual normalization. This bilinear observation is not an orthogonal
projection; transpose is intentional. Periods use cochain images and
remain a distinct topological control.
For boundary harmonic chain bases $\Phi,\Phi^\vee$, let $P,P^\vee$ be
periods of their images and $B_0=(\Phi^\vee)^{\mathsf T}M_{\rm own}\Phi$.
Square invertible period matrices give
\[
 G_{\rm per}=(P^\vee)^{-\mathsf T}B_0P^{-1}.
\]
The abbreviated $(P^\vee)^{-\mathsf T}P^{-1}$ requires $B_0=I$.
Zero-connection primal and dual/primal constructions agree only when they
describe the same period-normalized pairing; this does not equate historical
primal-frame and newer biorthogonal observation coordinates.
With primal/dual transports $M,M^\vee$, pairing preservation tests
$(M^\vee)^{\mathsf T}G_BM=G_A$. Using $M$ on both sides assumes coincident
transports in the declared trivialization.
This is distinct from positive quantum isometry. A complex bilinear form
is not a Krein inner product, which must be nondegenerate indefinite
Hermitian. Neither supplies an invariant positive sector, a
particle/antiparticle split or complete positivity automatically.

\section{Internal elimination and actual gluing}\label{sec:elimination}

For retained/internal coordinates $b,z$, invertible $K_{zz}(0)$ gives
\[
 z=-K_{zz}^{-1}K_{zb}b,\qquad
 S_G=K_{bb}-K_{bz}K_{zz}^{-1}K_{zb},\qquad S_Gb=0.
\]
This Schur complement is exact elimination, not a partial trace; internal
coordinates remain reconstructable. Singular blocks require compatibility
and null modes, not an unchecked pseudoinverse. The positive Hermitian
fixture $S=[-T,I]^\dagger W[-T,I]$, $W>0$, has zero energy exactly for $y=Tx$.
For two unit-weight constraints put $y=T_1x+u$, $r=z-T_2T_1x$.
Minimizing $\|u\|^2+\|r-T_2u\|^2$ gives
$u=(I+T_2^\dagger T_2)^{-1}T_2^\dagger r$ and energy
\[
 r^\dagger(I+T_2T_2^\dagger)^{-1}r .
\]
The zero relation composes to $z=T_2T_1x$. A target-built fixture is an
algebraic control, not scalar-geometric synthesis. Failure to synthesize
it does not invalidate the independent representation theorem. Actual composition must
identify interface cells/frames and compare the assembled glued relation.
Extra dimension does not prove arbitrary matrix realizability. A connected
positive scalar degree-zero connection graph has at most one zero mode:
$f_j=e^{i\theta_{ij}}f_i$ fixes all values from one vertex; inconsistent
holonomy forces zero. Higher-degree topology is a candidate, not a theorem.

\section{Quantum norms, tensor products and channels}\label{sec:quantum-contract}

Declare positive Hermitian $Q_{\rm in},Q_{\rm out}$ and test
$T^\dagger Q_{\rm out}T=Q_{\rm in}$. Orthonormal coordinates use
$\widehat T=Q_{\rm out}^{1/2}TQ_{\rm in}^{-1/2}$.
Only with a declared positive norm can a harmonic vector be normalized
as a quantum state. Stationary-action relaxation does not prove isometry.
Fix output-major vectorization:
\[
 |T\rangle\!\rangle=\sum_{i,j}\widehat T_{ij}|i\rangle_{\rm out}
 \otimes|j\rangle_{\rm ref},\qquad
 (I\otimes\langle\bar x|)|T\rangle\!\rangle=\widehat T x .
\]
The reference is the conjugate/dual input, not another physical boundary.
For nonzero $\widehat T$, normalize by $\|\widehat T\|_F$. Schmidt
coefficients are normalized singular values; the reference marginal is
$(\widehat T^\dagger\widehat T)^{\mathsf T}/\|\widehat T\|_F^2$, equal to
$I/d_{\rm in}$ for an isometry. This is operator-vector information,
not automatically physical-state entanglement.
For Kraus operators in orthonormal coordinates,
\[
 \mathcal E(\rho)=\sum_aK_a\rho K_a^\dagger,\qquad
 J_{\mathcal E}=\sum_a|K_a\rangle\!\rangle\langle\!\langle K_a| .
\]
CP is $J_{\mathcal E}\ge0$; TP is $\Tr_{\rm out}J_{\mathcal E}=I_{\rm in}$
in this unnormalized convention. Stinespring requires $V:A\to B\otimes E$,
$V^\dagger V=I$, and $\mathcal E(\rho)=\Tr_E(V\rho V^\dagger)$.
A third boundary in $B\oplus E$ is not that tensor product.

\subsection{Three distinct tensor constructions}
\begin{enumerate}
\item K\"unneth over a field:
$H^n(X\times Y)\cong\bigoplus_{p+q=n}H^p(X)\otimes H^q(Y)$.
An $n$-fold degree-one register occupies multidegree $(1,\ldots,1)$,
not generally all of $H^n$. Harmonic factorization needs a compatible
product complex/metric. K\"unneth is not limited to quantization level one.
\item Theta quantization gives $k^g$ sections from a rank-$2g$ lattice
with principal polarization and level line bundle; it is not $H^1$.
The modulus register uses $\tau$ to select a line in $H^1$; the theta
register uses $\tau$ as polarization and a section as state.
\item Ordinary root-system tensor multiplicities come from
Brauer--Klimyk; affine fusion comes from Verlinde/Kac--Walton.
Fusion is not the untruncated tensor product. Disjoint surfaces factor
state spaces; sewing/punctures require sums over labels and fusion spaces.
\end{enumerate}
Two level-two theta factors give four amplitudes, a $4\times4$ gate and
sixteen operator-vector amplitudes. A direct sum of two harmonic port
frames, or historical $2\times2$ selected-pair $\chi$, is not that operator.
Dimension matching alone does not define subsystems.

\section{Current representation acceptance contract}\label{sec:validation}

The present deliverable is a representation theorem with computational
checks of its hypotheses and operations, not a dynamics or synthesis
claim. Theorems hold exactly under their assumptions; numerical residuals
measure implementation accuracy. Tickets \#1121--\#1124 divide the work
as in Section~\ref{sec:next}. All four must identify which operations are
geometric computations and which are independent matrix oracles.
\begin{enumerate}[label=V\arabic*.,leftmargin=*]
\item State the classical/quantum type bridge, integral markings, positive
polarization, level, theta structure and coefficient/evaluation convention.
Use positive Hodge masses for the classical proof, or report a certified
alternative; do not import positivity from a complex bilinear pencil.
\item \#1121: represent $e_{p,q}$ by scalar translations; verify the crossing
law, involution, $d^2$ independent symbols, trace pairing, forward/inverse
maps and twisted convolution against independent matrices, including
non-Clifford and nonunitary operators. Preserve the pinned-boundary defect
as a distinct invariant control, not an optimization target.
\item \#1122: verify ordinary closed phase configurations, Wilson loops,
gauge/large-gauge invariance and neutral-pencil independence. Apply the
coordinate bridge and correct conjugate coherent action. Demonstrate
qubit coverage with several target rays including zeros/poles; keep
level-three coverage limitations and characteristic restrictions explicit.
\item \#1123: verify tensor-factor ordering, matrix units and Choi contraction
against composition, including rectangular maps, product controls, SWAP
and an entangling witness. Keep the wider-neck geometry study as a
separate representation/entropy experiment, not a requirement for the
algebra theorem or a claim about Ryu--Takayanagi.
\item \#1124: recover $A_2$ from scalar lattice lengths, its roots and dual
lattice; verify character dimensions, Weyl alternation, modular data and
fusion by both geometric alcove evaluation and independent folding.
The native fiber need not be synthesized.
\item Use fixed-seed held-out inputs distinct from fit samples, independent
formulas, conditioning/rank reports and truncation/precision convergence.
Report residuals before normalization/rounding and reject failed
hypotheses. Surface periods require symmetry, positive imaginary part,
and a justified complex structure; report $J^2+I$ and any normalization
rather than hiding it. Validate refinement where a discrete read
approximates continuum geometry.
\item Keep default \texttt{verify}/\texttt{theta} fixed-complex and read-only
with respect to engine evolution. Reuse existing Gram/pencil/theta
implementations. Preserve compatibility except for explicitly corrected
scientific expectations; add regressions demonstrating why they changed.
\item Each record declares its \texttt{certifies} scope: exact algebra
representation, sampled analytic-family check, approximate geometric
readout, or native realization. A native realization failure is not a
representation failure. No matrix connection, engine relaxation,
universality, speedup or entropy--area law is required or inferred.
\end{enumerate}

\section{Retained native implementation and acceptance contract}\label{sec:implementation}

These requirements preserve the earlier native-engine contract where that
path is exercised; they do not replace V1--V8 or require new synthesis.
Original IDs are retained; short verifier names are a separate namespace.
Existing incorrect scientific claims must be corrected with documented
regressions rather than preserved as acceptance criteria.

\subsection{Construction R1--R7}
\begin{enumerate}[label=R\arabic*.,leftmargin=*]
\item A native harmonic state is a whole cochain including boundary;
restrictions are readouts. A theta-register state is instead a section,
or its density symbol; these are different types.
\item Reuse Whitney pencils, \texttt{MultiCobordism}, fibers and shared
\texttt{emergence\_animation.py} runtime; keep \texttt{qubit\_animation.py} canonical.
\item Preserve pinned defaults and \texttt{--no-pin-boundary}; never remove boundary cells.
\item Preserve manifold/collar gates, stage-1 search and stage-2 scalar
length/phase relaxation. No matrix-valued connection or target in links,
measured frames or assembled Laplacians.
\item Targets may condition synthesis, never frozen recovery.
\item Preserve complex/Lorentzian geometry; do not silently impose quantum positivity.
\item Read the zero kernel, distinguish excited/generalized bands and
apply Section~\ref{sec:kernel-contract}.
\end{enumerate}

\subsection{Scientific readout S1--S7}
\begin{enumerate}[label=S\arabic*.,leftmargin=*]
\item Retain flat-torus inputs, markings, holomorphic fibers and own-state residuals.
\item Retain scalar collar synthesis and shared engine-unit schedule.
\item Record the actual objective: historical transfer/$\chi$ fitting
optimizes selected-state amplitudes, not the whole operator.
\item Recover whole-kernel relations each frame through periods and native
dual-frame whole-metric contractions using \texttt{PencilSchur.gramBlock}
and Section~\ref{sec:gram-contract}. The \texttt{gram} record specifies
\texttt{observation\_convention: dual\_frame\_whitney}. Older records lack
this field and used primal frames; normalizations need not agree even
at zero phase. Identity coordinate norms are explicit external encodings.
\item Preserve historical $\chi$, not as a full XY/two-qubit gate and
never as input to frozen recovery.
\item Publish \texttt{correspondence}: identifiability, ranks, conditioning,
basis/held-out errors, norm preservation, Choi data and a separate gate verdict.
\item Report the rank-two modulus host's obstruction to a four-dimensional
two-qubit gate despite objective convergence. Theta sphere joins give
product operators; entanglement needs a validated operator and product-input
witness, not merely a proposed circle seam.
\end{enumerate}

\subsection{Compatibility D1--D5}
\begin{enumerate}[label=D\arabic*.,leftmargin=*]
\item Preserve native collar/bridge gates and rollback.
\item Preserve own-surface frames and block-state residuals.
\item Keep bilinear pencil transfer distinct from whole transport and certified channels.
\item Keep readout in \texttt{qubit\_animation.py} with reusable helpers;
preserve CLI, geometry records, output handling, neutral runtime and
historical measurements while displaying actual whole-operator evidence.
\item \texttt{verify}/\texttt{theta} read fixed complexes seeded through
\texttt{seed\_qubit\_host}, without node, objective, readout optimization or
engine stages. \texttt{MultiCobordism.restriction} supplies one basis,
all marking periods, chain frame and certificate; \texttt{monodromy} adds
its fit. Preserve parameter/modulus-named records under
\path{~/cobordism-runs/functor-verify/} and
\path{~/cobordism-runs/theta-register/}, with a delimiting
\texttt{certifies} sentence. \texttt{run.sh} runs both.
\end{enumerate}

\subsection{Acceptance C1--C9}
\begin{enumerate}[label=C\arabic*.,leftmargin=*]
\item Live collar: kernel residual, input coverage, unique output and
basis/held-out whole-cochain witnesses; no target in recovery.
\item Independent complex, rectangular, hidden-mode, ambiguous-output,
incomplete-coverage, zero-map and rank-loss fixtures; rebase invariance.
\item Choi contractions, marginals, norms and composition;
no unitary certificate for nonunitary or zero maps.
\item Perturb scalar geometry: periods stay topological on fixed marked
collars, Gram observations may move. For pure gauge $g(v)=e^{i\phi(v)}$
fixed to one at port basepoints, both operators remain unchanged for real
and complex lengths. Otherwise the native convention gives
$T'=e^{-i(\phi_{\rm out}-\phi_{\rm in})}T$; normalized Choi marginals,
entropy and Schmidt coefficients remain invariant. Retain $\C^*$ dual
gauge-covariance controls.
\item Shared headless/live records, finite-safe JSON, panels naming
claims/obstructions and CLI/runtime regressions.
\item A failed gate realization remains failed when selected-state fitting converges.
\item Preserve \texttt{verify} at $10^{-12}$ on product/swap two-/four-torus
and identical-square hosts: rank, integrality, declared isotropy, flipped
control, closed and coclosed independent kernels and certificates before and after
jitter, fixed periods, exact representative/Gram motion, twist-group law,
direct sum and transport defect. Keep doubled-swap, $(1,10^{-5})$ exactness,
isometric-start, $10^{-7}$ rebase, same-convention zero-phase Gram and
$U(1)$/$\C^*$ gauge regressions; report counts and failures.
\item Preserve \texttt{theta} T1--T8: transformation law, norm, projective
law, generators, factorization, antiunitarity and controlled-$Z$.
Record level, truncation, fit/held-out samples, phases and factor ordering.
Algebraic generator controls do not prove native host synthesis.
\item Preserve the post-consolidation reads (Section~\ref{sec:evidence-merged}):
\texttt{verify --far-flips} F1--F4 and \texttt{theta --far-flips} T8/T9
(the Dehn twists and composition order on one host, grid $3$ for one pass,
grid $5$ for two); \texttt{theta --tori 4 --join tube} G1--G5 (the
genus-two far boundary: product monodromy, period matrix, coupling,
geometric-state entropy, neck sweep); \texttt{theta --surface decagon}
E1--E6 (the entangling monodromy of the order-five rotation, its fixed
period matrix, the Clifford gate, the product-collar control). Report
counts and failures; records under \texttt{\string~/cobordism-runs/}.
\end{enumerate}

\section{Root representation and optional native synthesis}\label{sec:root-program}

The root-lattice and fusion representations are current objectives;
native fiber synthesis is optional. No arbitrary matrix connection is
introduced. The internal coroot lattice is independent of boundary
cohomology: boundary mapping classes act on $H^1(\Sigma;\Z)$, whereas
the Weyl group acts on $Q^\vee$. Even when the abstract groups coincide,
these are not the same action.

For a torus boundary the abelianized Cartan phase lattice is
$H^1(\Sigma;\Z)\otimes Q^\vee$, with the product of the surface intersection
form and the positive internal root form, followed by the declared affine
reduction. The surface supplies $\tau$; internal scalar lattice geometry
can supply $A_2$, its roots, weights and reflections. The series notation
$\Omega=\tau\operatorname{Gram}(Q^\vee)$ must retain the nonprincipal
polarization and characteristic set of Section~\ref{sec:root-coordinates}.
A genus-$r$ boundary whose principal period matrix happens to equal
$\tau A$ is not, by that equality alone, the root register.

The present acceptance claim is that the lattice/alcove construction
reproduces character and fusion algebra. Later native synthesis would
add a sector map, certified whole-kernel readouts and actual sewing;
braiding would add labeled punctures and compatible associativity/braiding
maps. These are distinct extensions, not missing prerequisites for the
current representation result. Extra dimensions permit candidate encodings
but do not themselves prove native realizability.

\section{Neck geometry and encoded-state entropy}\label{sec:neck-entropy}

\begin{definition}[Entropy in a declared tensor register]
For a normalized pure section coefficient vector in
$\mathcal Q_1\otimes\mathcal Q_2$, let $p_i$ be its squared Schmidt
coefficients. Its reduced von Neumann entropy in bits is
$S=-\sum_{p_i>0}p_i\log_2p_i$. This measures the encoded state relative to
the declared factors, not an intrinsic function of one period entry.
\end{definition}

\begin{proposition}[Weak-coupling entropy is generally not quadratic]\label{prop:entropy-asymptotic}
Suppose a two-qubit normalized state varies analytically with a real
parameter $\epsilon$, is product at zero, and has a nonzero first-order
component normal to the product-state manifold. Its small Schmidt
probability obeys $p=C\epsilon^2+O(\epsilon^3)$, $C>0$, and
\[
 S=-C\epsilon^2\log_2(C\epsilon^2)
       +\frac{C\epsilon^2}{\ln2}+o(\epsilon^2).
\]
For level-two theta states with
$\Omega=\operatorname{diag}(\tau_1,\tau_2)
+\epsilon\begin{psmallmatrix}0&1\\1&0\end{psmallmatrix}$
at $z=0$, the first derivative vanishes. If the second-order entangling
component is nonzero, $p\sim C\epsilon^4$ instead.
\end{proposition}
\begin{proof}
The small singular value is first order in the normal component, so its
square is $C\epsilon^2+O(\epsilon^3)$. Expand the binary entropy
$-p\log_2p-(1-p)\log_2(1-p)$.
For the theta family,
$\partial_\epsilon\theta|_0=(2\pi i k)^{-1}
\theta'_{j_1}(z_1;\tau_1)\theta'_{j_2}(z_2;\tau_2)$.
Both level-two sections are even functions of their argument, so their
first derivatives vanish at zero. The leading entangling amplitude can
then be second order; special parameters may cancel it as well.
\end{proof}

The star-tube experiment retains wider disk attachment, waist/length
sweeps and larger grids. Treat $|\Omega_{12}|\sim0.3$--$0.5$ as an
exploration target, not a promised output. Record full $\Omega$, markings,
diagonal drift, period-reader errors, mesh resolution and the same fixed
real holonomy samples $(a,b)$ at each point; convert to $z=b-\Omega a$.
Synthetic asymptotic controls at fixed $z,\tau_1,\tau_2$ are different
from geometry sweeps in which these quantities move. Report both Schmidt
probabilities and entropy; fit the appropriate logarithmic asymptotics,
separating $z=0$ from generic points.
Topology/positivity at one strong point does not establish convergence.
A physical degeneration must be checked on actual geometry; manually
zeroing $\Omega_{12}$ is only a factorization control.

This is encoded-state entanglement, not a Ryu--Takayanagi validation.
No independent extremal-area observable, area normalization or holographic
dictionary is supplied. Stronger neck coupling is not required to prove
the full operator algebra, and does not by itself realize a root fiber.

\section{Information and geometric diagnostics: unresolved definitions}

The earlier follow-up (\#1095 in the pinned spec) asks for total-system
entropy against engine units and changing curvature/charge heatmaps.
These are not operator-vector entropy. First declare a positive normalized
joint $\rho$, subsystem factors and evolution/coarse-graining rule. Record
$S(\rho)=-\Tr(\rho\log\rho)$, reduced entropies and
$I(A:B)=S(\rho_A)+S(\rho_B)-S(\rho_{AB})$, with log base and engine schedule.
A globally pure closed state has zero total entropy while local entropies
may change; their sum is not total entropy.
Mutual information measures total correlations, not an intersection of
directional flows or entanglement alone. An MI-dependent inner product at
the $3$-simplex needs a joint-state construction and a Hermitian
positive-semidefinite kernel for every finite set; symmetry/nonnegativity
alone are insufficient. Nondegeneracy is additionally needed.

Unobserved harmonic remainder is not automatically an environment or lost
entanglement. Retain it in reconstruction until projection, conditioning or
partial trace is declared with normalization and CP/TP status.
Transport/merging must propagate the joint state and correlations under
that operation. Curvature heatmaps need an estimator, dual-cell support,
Lorentzian/complex convention, units and normalization; separate
\emph{mass curvature} needs a defined mass observable and geometric relation.
Charge maps need an observable, sign, support and conservation test.
No identity equating curvature, charge, entropy or remainder is proved.

\part{Evidence, reconciliation and references}

\section{What the code measures}

The pinned specification reports that \texttt{qubit\_animation.py verify}
reads a seeded collar at tolerance $10^{-12}$: rank and integrality of $M$ (R1--R4);
isotropy for the declared form on the orthonormalized restriction subspace
with the flipped control nonvanishing (R5, Remark~\ref{rem:sign} and P4);
the harmonic certificates (H1--H3, the caveat after Theorem~\ref{thm:hodge});
metric independence of $M$ against an $O(1)$ motion of the representative
that is exactly a coboundary (J1--J3, Theorem~\ref{thm:invariance}); the
twist-group composition (C1, P11); the four-torus direct sum (D1,
Proposition~\ref{prop:sphere}); and the covariant transport defect (U1,
Section~\ref{sec:metric} and Remark~\ref{rem:gauge}). The subcommand
\texttt{theta} fits $\rho_2$ for the generators and the seeded monodromies
(T1--T8, Sections~\ref{sec:weil}--\ref{sec:seam}). These are source-reported
fixture results, not new runs or proofs by numerical precision.
The reported small residuals lie at $10^{-16}$--$10^{-14}$; metric-dependent
defects can be order one or small for matching boundaries.

\begin{center}
\begin{tabular}{@{}>{\raggedright\arraybackslash}p{0.14\textwidth}>{\raggedright\arraybackslash}p{0.79\textwidth}@{}}
\toprule
status & statement and limitation \\
\midrule
theorem & $L_W$ is Lagrangian for oriented manifold boundaries over a field;
harmonic readout additionally needs a certified cohomology identification.\\
conditional & $M$ is integral unimodular under integral lattice isomorphisms,
not merely invertible period matrices (Theorem~\ref{thm:integer}).\\
theorem & Relations compose by cohomological gluing; disjoint union gives
direct sum. Symbol contraction is a separate exact composition theorem;
native seam gluing requires independently assembled complexes.\\
theorem & Theta section spaces factor; even-level Weil operators are
Clifford. This is a new register, not equality with harmonic $H^1$.\\
reported & Bilinear transport defect $0.46$--$0.58$ generically, about
$10^{-14}$ for identical tori; conventions are in Section~\ref{sec:gram-contract}.\\
measured (Section~\ref{sec:evidence-merged}) & Gluing along the far surface
of one host by layered flip passes, with the Dehn twists $T_A$, $T_B$, $S$
and their non-commuting composition; a genus-two far boundary whose period
matrix couples the tori while the monodromy stays a product; an entangling
monodromy from the order-five rotation of a symmetric genus-two surface.\\
current & Geometric operator algebra (\#1121), holonomy/coherent encoding
(\#1122), tensor/Choi composition and neck entropy (\#1123), root-lattice
characters and fusion (\#1124); see V1--V8.\\
optional & Native seam and through-neck gate synthesis, positive physical
dynamics and non-abelian fiber/sewing implementation. These do not gate
the representation result. Complex bilinear pairings remain distinct
from positive quantum norms.\\
\bottomrule
\end{tabular}
\end{center}

\subsection{Numerical evidence retained from the specification}
For grid $3$, seed $7$ and tolerance $10^{-12}$, the source reports:
\begin{itemize}
\item Period integrality residual at most $10^{-14}$; product/swap maps
equal identity/swap; declared isotropy at most $1.1\,10^{-14}$, versus
flipped controls $1.41$ (two tori) and $2.00$ (four).
\item A $75\%$ complex length jitter changes $M$ by at most $1.6\,10^{-14}$
but moves the canonical representative by $0.49$--$0.77$, with exact
difference residual at most $4\,10^{-14}$. Reported Gram motion is
$0.18$--$0.45$.
\item Twist-group composition residual at most $1.5\,10^{-14}$; sphere-join
off-block norm at most $5\,10^{-15}$. The marking maps on the current
square triangulation commute, so this does not test composition order.
\item Closed/coclosed residuals at most $2.5\,10^{-14}$, nullity/rank
$2$ or $4$, independent span residual at most $1.2\,10^{-14}$ and
projector idempotency at most $7\,10^{-16}$.
\item Level-two theta-law residuals at most $6.8\,10^{-16}$;
unitarity $1.3\,10^{-15}$; projective group law $9.5\,10^{-16}$;
Hadamard/phase checks $4.4\,10^{-16}$; factorization $1.9\,10^{-16}$;
Kronecker law $6.2\,10^{-16}$; controlled-$Z$ $5.3\,10^{-16}$;
polarization independence $8.8\,10^{-16}$. Product/CZ operator Schmidt
ranks are one/two. The algebraic CZ control is not a native entangling host.
\item Reported \texttt{verify} counts: $17/17$ for each of product, swap
and identical-square two-torus fixtures; $27/27$ for each four-torus
product/swap fixture. Reported \texttt{theta} counts:
$8/8,8/8,10/10,10/10$ on two-/four-torus product/swap fixtures.
\end{itemize}
No engine dynamics, actual interior-seam composition or scalar-geometric
non-abelian register was established by these fixed-complex reads.
A topology search move need not change the marked relation: a
triangulation-preserving equivalence can leave it unchanged.

\subsection{Measurements after consolidation}\label{sec:evidence-merged}
Made on 14 September 2026 on the fixed complexes named, tolerance
$10^{-12}$, grid $3$ and seed $7$ unless stated; records under
\path{~/cobordism-runs/functor-verify/} and
\path{~/cobordism-runs/theta-register/}; merged as Tessera
\#1116, \#1119 and \#1120.
\begin{itemize}
\item \textbf{Tube-joined collars} (\#1116; \texttt{theta --tori 4 --join tube}).
Two collars joined by a prism tube between one far triangle of each: the far
boundary is one genus-two surface, $\partial W = T^2\sqcup T^2\sqcup\Sigma_2$,
Betti numbers $[1,4,2,0]$, restriction rank $4 = b_1(\partial W)/2$, boundary
Euler characteristics $[-2,0,0]$, declared isotropy $4.6\,10^{-15}$ against
the flipped control $2.0$ (G1). The monodromy from the near tori to
$\Sigma_2$ is $I_4$ (product collars) or swap$\,\oplus\,$swap, off-block
norm $2.6\,10^{-15}$, operator Schmidt rank one (G2): the tube alone couples
no cycle. The surface's own Riemann period matrix, read from its harmonic
$1$-forms with the torus reader's conventions at genus two, has
$\operatorname{Im}\Omega$ positive definite (eigenvalues $0.83$, $1.08$),
$\Omega_{12} = 1.45\,10^{-4} - 1.88\,10^{-4}i$ on the default neck and is
symmetric to $2.4\,10^{-15}$ at every neck read (G3). The geometric
(coherent) state at a holonomy point, in the basis the identity monodromy
identifies with $\Theta_2\otimes\Theta_2$, has entanglement entropy exactly
$0$ for the diagonal control ($5\,10^{-30}$) and $4\,10^{-7}$ bits mean on
the default neck (G4). Over $25$ necks (waist $0.25$--$4$, length
$0.25$--$4$) $|\Omega_{12}|$ runs $5\,10^{-8}$ to $2.4\,10^{-2}$, falling
approximately exponentially with tube length in the reported sweep.
The reported entropy range is $5\,10^{-14}$ to $2.4\,10^{-3}$ bits mean;
its historical description as proportional to $|\Omega_{12}|^2$ was
a finite-range observation, not the asymptotic theorem.
The diagonal of $\Omega$
approaches the far tori's moduli as the neck pinches ($0.099\to0.037$): the
separating degeneration (G5). Counts $11/11$, $11/11$ (swap), $12/12$
(sweep). This is a weak-coupling encoded-state observation, not an
entropy--area law. Historical averaged values used the then-current
sampling procedure and are not paired-sample convergence evidence.
The corrected asymptotics and controls are in
Section~\ref{sec:neck-entropy}.
\item \textbf{Dehn twists by layered flips} (\#1119;
\texttt{verify --far-flips}, \texttt{theta --far-flips}). One tetrahedron per
flipped boundary edge of the far torus (nine per pass on grid $3$): rows
for $T_B$, columns for $T_A$, diagonals then the quarter turn for $S$; the
flipped grid is the grid in a sheared basis, relabelled by the lattice map
$L$, and the far marking reads $M = L^{\mathsf T}$. Measured
$T_B = \begin{psmallmatrix}1&1\\0&1\end{psmallmatrix}$,
$T_A = \begin{psmallmatrix}1&0\\1&1\end{psmallmatrix}$,
$S = \begin{psmallmatrix}0&-1\\1&0\end{psmallmatrix}$ with residual
$\le 4\,10^{-15}$, of infinite order for the twists (F2); the flipped far
torus's own modulus in the relabelled marking equals
$(M_{21}+M_{22}\tau)/(M_{11}+M_{12}\tau)$ to $8\,10^{-16}$ (F3);
$\rho_2 = \operatorname{diag}(1,i)$, $H\operatorname{diag}(1,-i)H$, $H$ at
$\le 5\,10^{-16}$ (T8). On grid $5$: $T_B^2$, $T_AT_B$ and $T_BT_A$ read as
the products of their passes, the two orders differing by one (F4) and
their Weil matrices at projective distance $1.0$ (T9); with the swap,
$M = M(\text{flips})\cdot\text{swap}$. Counts $20/20$ (one pass),
$21/21$ (two passes), $8/8$, $9/9$. On grid $3$ a second pass collides with
an existing edge and is refused by name.
\item \textbf{The entangling collar} (\#1120; \texttt{theta --surface decagon}).
The decagon with opposite sides identified, coned and barycentrically
subdivided ($28$ vertices, $90$ edges, $60$ faces, $\chi = -2$), the flat
decagon metric, its rotations as vertex permutations (simplicial isometries
of order ten), a symplectic marking by integer symplectic reduction of the
intersection form read from the surface (cup pairing of the harmonic basis,
sign calibrated on the flat torus). The collar relabelled by $r^2$ reads
$M = \begin{psmallmatrix}-1&1&0&0\\-1&0&1&1\\0&0&0&1\\-1&0&0&0\end{psmallmatrix}$
on $(A_1,A_2,B_1,B_2)$ as predicted from the side classes, $\det 1$,
$M^{\mathsf T}JM = J$, $M^5 = I$, mixing (E2); the surface's own period
matrix is fixed by $M$ to $1.5\,10^{-15}$, $\operatorname{Im}\Omega$ with
eigenvalues $0.46$, $1.66$ (E3); $\rho_2(M)$ unitary (residual $6\,10^{-16}$),
$\rho^5\propto I$ ($1.5\,10^{-15}$), operator Schmidt rank four (E4); $16$
of the $36$ product stabilizer states are sent to Bell pairs (exactly one
bit), projective distance $1.00$ to every $A\otimes B$ and $0.77$ to every
$\mathrm{SWAP}\cdot A\otimes B$ over the $24\times24$ single-qubit Cliffords
(E5); the product collar reads $M = I$ with entropy $5\,10^{-29}$ (E6).
$r^4$: order five, Schmidt rank two, entangling; $r$: order ten, entangling;
$r^5 = -I$: Schmidt rank one, no entangling power. Counts $13/13$ each.
Reduced mod $2$ an element of order five of $\Sp(4,\mathbb F_2)\cong S_6$
lies in no conjugate of $S_3\wr\Z/2$, so no such gate is a local Clifford
or a swap times one: the entangling capacity is forced by the order.
\end{itemize}

\paragraph{Historical modular conventions.}
The reported theta matrices above used the historical direct-monodromy
input convention. They certify those recorded fits and witnesses, not
automatically covariance with the simultaneously reported boundary modulus.
Current tests must apply Proposition~\ref{prop:coordinate-bridge} and
Proposition~\ref{prop:coherent-covariance}, explicitly distinguishing raw
period matrices, standard Siegel matrices, and value/state actions.
Likewise a finite-dimensional discrete period read is not promoted to an
exact Riemann-surface theorem by symmetry and positivity alone.

\subsection{What is measured next}\label{sec:next}
Owner-approved representation objective, revised 14 September 2026.
Implementation tickets remain open; this document and their specifications
are revised in the documentation PR for \#1131.
\begin{enumerate}
\item \#1121: complete geometric operator algebra. Implement reduced
cycle/translation symbols, scalar crossing phases, involution, trace and
twisted convolution. Test the $d^2$ basis, round trips for arbitrary
operators, adjoints and multiplication against independent matrices.
Retain fixed-boundary pairing invariance as a distinct control; do not
add relaxation to the fixed-complex verifier or expect its defect to fall.
\item \#1122: holonomy-defined states. Keep the neutral harmonic pencil,
construct ordinary closed phase configurations and read Wilson loops
independently. Fix $z=b-\Omega a$, the period/modular bridge, conjugate
coherent action and antiunitary factorization. Check complete qubit
encoding and its non-unique inverse. At level three use a named
characteristic-preserving subgroup or explicitly track characteristics;
do not claim coherent coverage of every qutrit ray.
\item \#1123: tensor and Choi algebra plus the retained neck study.
Verify symbol tensor products, dual-factor contractions, product/SWAP/CZ
controls, and state/channel distinctions. Retain the star disk and larger
grid sweeps; fix sample pairing, period accuracy/refinement and logarithmic
entropy asymptotics. A strong neck is an exploration result, neither an
RT law nor a prerequisite for representing entangling operators.
\item \#1124: root-lattice representation. Derive $A_2$ roots/reflections
and the weight lattice from scalar internal lattice geometry, then use
the correct nonprincipal bundle and $P/KQ^\vee$ characteristic set.
Fit normalized Weyl--Kac characters against Kac--Peterson on held-out
points. Compare fusion by geometric alcove evaluation, Verlinde and
independent Kac--Walton folding; include $A_1$ and the correct even-lattice
Frenkel--Kac controls. Host-derived $\tau$ is an input source, not a demand
to synthesize the root fiber.
\end{enumerate}
Common requirements are V1--V8. A native seam between separately seeded
cobordisms, a through-neck Dehn twist, splitting back to separate boundary
tori, physical dynamics and braid synthesis are optional future work.
Their absence does not weaken the stated algebra-isomorphism theorem.

\section{Source coverage and resolved conflicts}\label{sec:reconciliation}

\begin{longtable}{@{}>{\raggedright\arraybackslash}p{0.29\textwidth}>{\raggedright\arraybackslash}p{0.64\textwidth}@{}}
\toprule
source material & maintained location or correction\\
\midrule
\endhead
Original derivation 0--5 &
Sections~\ref{sec:cochains}--\ref{sec:composition}: harmonic/cohomology bridge,
mutual annihilators, graph condition, dual-input vectorization and
class-level gluing. Literal space/vector equality and the unqualified
Atiyah TQFT claim are withdrawn.\\
Original derivation 6--9 &
Sections~\ref{sec:metric}, \ref{sec:kernel-contract}--\ref{sec:quantum-contract}:
metric effects, quantum norms, channels and twisted cohomology.
Stationary action does not prove isometry; a generic complex bilinear
form is not Krein. The earlier blanket exclusion of nondegenerate
length-cochain states is not proved here. Determinant weighting alone
does not establish the operator correspondence.\\
Spec scientific question and R1--R7 &
Sections~\ref{sec:kernel-contract}, \ref{sec:implementation};
whole scalar geometry, target-free recovery and separate scientific verdicts.\\
Spec mathematical contracts &
Sections~\ref{sec:kernel-contract}--\ref{sec:quantum-contract} plus
Sections~\ref{sec:restriction}--\ref{sec:seam}; positive Hodge identity,
torus encoding, identifiability, Schur elimination, Choi/channel conventions,
harmonic certificates and theta construction.\\
Spec S1--S7, D1--D5, C1--C9 &
Section~\ref{sec:implementation} retains the native identifiers and
compatibility requirements; V1--V8 govern the current representation
objective. Historical residuals and counts retain their original meaning.\\
Prior reading, theta and seams &
Sections~\ref{sec:weights}--\ref{sec:clifford}: polarization/line-bundle data,
missing modular determinant factor, even-level/lift caveats, antiunitarity
and product-input entanglement witness. A circle overlap or mixing matrix
alone is not an entangling-realization theorem.\\
Prior reading, roots and fusion &
Sections~\ref{sec:lie}--\ref{sec:density}, \ref{sec:root-program}:
roots, coroots, weights, characters, Brauer--Klimyk, affine level/alcove,
Weyl--Kac, Sugawara, Kac--Peterson, Verlinde, Kac--Walton,
Frenkel--Kac--Segal and qualified density. Root rank is not forced by
the boundary's $H^1$ rank; Weyl isotypic sectors are not alcove labels.\\
Tensor and level terminology &
Section~\ref{sec:quantum-contract}: K\"unneth, theta factorization,
ordinary representation tensor products and fusion are distinct.
The level-one equivalence is the simply-laced lattice construction,
with $\mathfrak{su}(2)_1$ corresponding to abelian theta level $2$.\\
Representation refinement &
Sections~\ref{sec:phase-space}, \ref{sec:operator-algebra},
\ref{sec:coherent}, \ref{sec:root-coordinates}, \ref{sec:validation}:
classical/quantum types, complete geometric algebra, covariance and qubit
coverage, root polarization, and current ticket acceptance. Native synthesis
is explicitly not required.\\
Information follow-up &
The preceding diagnostics section retains entropy/engine-unit,
curvature and charge objectives, with unresolved state/observable
definitions explicit rather than claimed measurements.\\
\bottomrule
\end{longtable}
This consolidation changes documentation only. Companion maps, problems
and solutions are retained as historical study aids; their numbering is
referenced where useful, but conflicting claims defer to this source.

\section{Where this sits}
\begingroup\small
\begin{itemize}[nosep]
  \item Gelca--Uribe,
    \href{https://arxiv.org/abs/1006.3252}{From classical theta functions to
    topological quantum field theory}: scalar translation/crossing algebra,
    theta quantization and the linearized decorated-cobordism interpretation.
  \item Manoliu,
    \href{https://arxiv.org/abs/dg-ga/9610001}{Abelian Chern--Simons theory}:
    a scalar-$U(1)$ TQFT completion with extended-manifold gluing data.
    It is not derived from the engine's Regge relaxation.
  \item Riemann--Roch for elliptic curves (e.g.\ Hartshorne,
    \emph{Algebraic Geometry}, Ch.~IV): the degree-two line-bundle argument
    in Theorem~\ref{thm:qubit-coverage}.
  \item Ryu--Takayanagi,
    \href{https://arxiv.org/abs/hep-th/0603001}{Holographic derivation of
    entanglement entropy from AdS/CFT}: an extremal-area correspondence
    with additional physical hypotheses, not a name for a theta neck curve.
  \item Atiyah, \emph{Topological quantum field theories} (1988): the
    tensor-monoidal axioms against which the direct-sum relation functor
    of Section~\ref{sec:composition} must be distinguished.
  \item Mumford, \emph{Tata Lectures on Theta I}: Definition~\ref{def:theta},
    the transformation law, and the Heisenberg group of Section~\ref{sec:theta}.
  \item Weil, \emph{Sur certains groupes d'op\'erateurs unitaires} (1964):
    the representation of Section~\ref{sec:weil}.
  \item Gottesman, \emph{The Heisenberg representation of quantum computers}
    (1998): why Clifford is classically simulable.
  \item Humphreys, \emph{Introduction to Lie Algebras and Representation
    Theory}, Ch.~II--III and VI: root systems, the Weyl group and the
    character formula of Sections~\ref{sec:lie}--\ref{sec:weyl}.
  \item Fulton--Harris, \emph{Representation Theory}, Lectures~11--13:
    $\mathfrak{sl}(2)$ and $\mathfrak{sl}(3)$ worked by hand.
  \item Kac, \emph{Infinite dimensional Lie algebras}, Ch.~6 (the affine
    Weyl group), 10--13 (integrable modules, the character formula and the
    modular matrices of Sections~\ref{sec:affine}--\ref{sec:verlinde}).
  \item Frenkel--Kac, \emph{Basic representations of affine Lie algebras
    and dual resonance models} (1980): Theorem~\ref{thm:fk}.
  \item Goddard--Olive, \emph{Kac--Moody and Virasoro algebras in relation
    to quantum physics} (1986): the Sugawara construction.
  \item Di Francesco, Mathieu, S\'en\'echal, \emph{Conformal Field Theory},
    Ch.~14--16: the $\mathfrak{su}(2)_k$ formulas, Verlinde, Kac--Walton.
  \item Ng--Schauenburg, \emph{Congruence subgroups and generalized
    Frobenius--Schur indicators} (2010):
    \href{https://arxiv.org/abs/0806.2493}{congruence and finiteness on the torus}.
  \item Freedman--Larsen--Wang, \emph{The two-eigenvalue problem and
    density of Jones representation of braid groups} (2002):
    \href{https://arxiv.org/abs/math/0103200}{density with representation-specific hypotheses}.
  \item NIST DLMF, \href{https://dlmf.nist.gov/21.5}{Section 21.5,
    modular transformations}: the determinant factor and characteristics.
  \item Arnold--Falk--Winther,
    \href{https://arxiv.org/abs/0906.4325}{Finite element exterior calculus}:
    positive Hodge theory, not positivity of a complex Lorentzian pencil.
  \item Watrous,
    \href{https://cs.uwaterloo.ca/~watrous/TQI-notes/TQI-notes.05.pdf}{Channel representations}:
    Choi, Kraus and vectorization conventions.
  \item Pinned implementation context:\\
    \path{simplicial_qubit_spec.md}\\
    \path{include/chainhodge/ChainHodge.h}\\
    \path{include/cobordism/PencilLayer.h}\\
    \path{include/cobordism/MultiCobordism.h}\\
    \path{tessera/quantum/theta_register.py}.
\end{itemize}
\endgroup

\end{document}

